\chapter{Bott--Chern classes and singular loci}
\label{chap:bott-chern-loci}

A closed class determines loci where its positive currents must be singular. Allowing an arbitrarily small negative error gives the non-nef locus; requiring strict positivity gives the non-K\"ahler locus. A Gauduchon trace makes the relevant current families compact, while analytic regularization turns these definitions into geometric sets. On manifolds admitting a K\"ahler modification, finite restricted masses give a numerical test for bigness after restriction. For a divisorial trace with bounded remainder and nef residual class, the residual intersection computes the trace mass; its positivity supplies a strict current by the diagonal construction.

The same trace pairing constrains divisorial components of Monge--Amp\`ere subsolutions at their maximal coefficient. An elliptic threefold makes the closed cones, section growth and degeneration of the variational volume explicit. On the ample locus of a big class on a K\"ahler manifold, a smooth positive volume form gives a smooth solution metric. Its completeness depends on the behavior near the omitted set. A blowup example gives finite distance to that set; a determinant estimate in the normal direction gives a sufficient completeness criterion and its sharp logarithmic exponent.

\section{Closed classes with Hermitian lower bounds}\label{sec:bc-hermitian-bounds}

Testing a current against a Gauduchon form eliminates its potential. The resulting fixed trace bounds all almost positive representatives of a class.

Fix a Hermitian metric $\omega$ on $X$. The ordinary real Bott--Chern group is
\[
 H^{1,1}_{\mathrm{BC}}(X,\R)
 =\frac{\{\theta\in C^\infty(X,\Lambda^{1,1}_{\R}):
                                      \mathrm{d}\theta=0\}}
        {\ddc C^\infty(X,\R)}.
\]
Fix a class $\alpha\in H^{1,1}_{\mathrm{BC}}(X,\R)$ and a smooth closed real $(1,1)$-representative $\theta$. Representatives that are currents have the form $T=\theta+\ddc u$; an almost positive such current has a quasi-plurisubharmonic potential $u$.

\Needspace{11\baselineskip}
\begin{definition}\label{def:bc-class-positivity}
The class $\alpha$ is \emph{pseudoeffective} if it contains a positive current. It is \emph{nef} if, for every $\varepsilon>0$, there is a smooth real $u_\varepsilon$ with
\[
 \theta+\ddc u_\varepsilon\geq-\varepsilon\omega.
\]
It is \emph{big} if it contains a closed current $T\geq\delta\omega$ for some $\delta>0$. Such a current is called a \emph{K\"ahler current}.
\end{definition}

These conditions are independent of the representative and the auxiliary metric. Changing the representative adds a smooth $\ddc$-term to the potential, and any two Hermitian metrics on a compact manifold bound each other by constant multiples.

For a pseudoeffective $\alpha$, write
\[
 \alpha[-\varepsilon\omega]
 =\{T\in\alpha:T\geq-\varepsilon\omega\}.
\]
\begin{proposition}\label{prop:bc-gauduchon-compactness}
For each fixed $\varepsilon>0$, the currents in $\alpha[-\varepsilon\omega]$ have uniformly bounded distributional mass.
\end{proposition}
\begin{proof}
Choose a Gauduchon metric $\gamma$. For $T=\theta+\ddc u$ in this family, the current $T+\varepsilon\omega$ is positive and
\[
 \int_X(T+\varepsilon\omega)\wedge\gamma^{n-1}
 =\int_X(\theta+\varepsilon\omega)\wedge\gamma^{n-1}.
\]
The right side is fixed. Comparison of smooth Hermitian metrics converts this trace into a bound for all coefficients of the positive current $T+\varepsilon\omega$. Subtracting the fixed smooth form gives the asserted bound for $T$. In dimension one the same argument uses Stokes directly.
\end{proof}

\begin{proposition}\label{prop:psef-cone-compact-base}
The pseudoeffective cone in $H^{1,1}_{\mathrm{BC}}(X,\mathbb R)$ is closed and convex and contains no nonzero line. For a fixed Gauduchon metric $\gamma$, the functional
\[
 \ell_\gamma(\alpha)=\int_X\alpha\wedge\gamma^{n-1}
\]
is strictly positive on its nonzero elements, and the slice of pseudoeffective classes with $\ell_\gamma=1$ is compact. This slice may be empty. Pullback and pushforward by a compact modification are continuous linear maps on Bott--Chern cohomology and preserve pseudoeffectivity.
\end{proposition}
\begin{proof}
Use the finite-dimensional smooth representative space of \cref{lem:constant-trace-bc-representative} to give Bott--Chern cohomology its usual finite-dimensional topology. Any choice of a basis or of the Gauduchon metric gives the same topology. The functional is well-defined by current Stokes and continuous by smooth representatives. If $T\geq0$ represents $\alpha$, then $\ell_\gamma(\alpha)$ is its trace mass. It is nonnegative, and equality forces $T=0$, hence $\alpha=0$. Convexity follows by adding positive currents. If $\alpha$ and $-\alpha$ are both pseudoeffective, their nonnegative masses add to zero; both classes are zero.

Let pseudoeffective classes $\alpha_j$ converge to $\alpha$, and choose positive closed representatives $T_j$. Their masses $\ell_\gamma(\alpha_j)$ are bounded. Weak compactness gives a subsequence converging to a positive closed current $T$. Writing $B_\gamma$ for the continuous smooth representative map of \cref{lem:constant-trace-bc-representative}, continuity of $B_\gamma(T_j)$ shows that the class of $T$ is $\alpha$. Thus the cone is closed. For its mass-one slice, \eqref{eq:constant-trace-representative-bound} bounds $B_\gamma(T)$ in the finite-dimensional representative space. The slice is closed and bounded there, hence compact.

The class maps of a modification are linear and well-defined by the smooth-to-current description of cohomology, so they are continuous between finite-dimensional spaces. Pullback preserves positive closed currents by local psh potentials, as in \cref{prop:closed-current-modification-continuity}; proper pushforward preserves positivity and closedness. These representatives prove the assertions about the cones.
\end{proof}

Compactness of this linear slice does not give a bound for Monge--Amp\`ere masses of potentials. The functional tests one closed current against a fixed smooth form; it is not a top-degree volume of the class.


This is also why a nef class is pseudoeffective: take smooth representatives with errors tending to zero and a weakly convergent subsequence. Closedness persists in the limit. The continuous representatives of \cref{lem:constant-trace-bc-representative} show directly that the Bott--Chern class also persists. Equivalently, normalized potentials have an $L^1$-convergent subsequence by the compactness of quasi-plurisubharmonic functions. The limit current is positive.

\subsection{Divisorial components of threshold subsolutions}

The linear Gauduchon threshold of \cref{cor:gauduchon-threshold} controls how much a divisorial component lowers the coefficient in a Monge--Amp\`ere inequality. For this subsection, fix a Hermitian form $\omega$ on a connected compact complex manifold of dimension $n\geq2$, a smooth probability volume $\mathrm{d}V$, and a nonzero $0\leq f\in L^1(\mathrm{d}V)$. Write $p=f^{1/n}$ and $b=b_\omega(f)$.

\begin{proposition}\label{prop:threshold-divisorial-constraint}
Suppose $u\in\PSH(X,\omega)$ satisfies
\[
 \mathrm{MA}^{\mathrm{ext}}_\omega(u)\geq c f\,\mathrm{d}V,
 \qquad \omega+\ddc u\geq[E],
\]
where $c>0$ and $E$ is an effective real divisor. Let $\beta_E$ be a smooth closed representative of its Bott--Chern class, and put $\kappa=\kappa_\omega(\beta_E)$. Then $b>0$, $0\leq\kappa<1$ and
\begin{equation}\label{eq:divisorial-coefficient-loss}
 c\leq b(1-\kappa)^n.
\end{equation}
If $c=b$, every smooth maximizing sequence $(T_j)$ for $D_p$ satisfies
\[
 \int_X\beta_E\wedge T_j=\int_X[E]\wedge T_j\longrightarrow0.
\]
Every weak cluster current $T$ therefore satisfies $\int_X\beta_E\wedge T=0$, and the class of $E$ is not big. If also $p\in\mathcal D_\omega^0$, each such cluster current maximizes $D_p$.
\end{proposition}
\begin{proof}
The subsolution characterization gives $b\geq c>0$. For a smooth strictly positive normalized Gauduchon form $T$, \cref{lem:unbounded-gauduchon-trace-density} shows that
\[
 R=(\omega+\ddc u)\wedge T-c^{1/n}p k_T\,\mathrm{d}V
\]
is positive, with mass $1-c^{1/n}D_p(T)$. The density term gives no mass to the support of $E$, whereas $[E]\wedge T$ is supported there. Restricting the positive measure $R$ to that support and using $\omega+\ddc u\geq[E]$ gives
\[
 \int_X[E]\wedge T\leq R(X)=1-c^{1/n}D_p(T).
\]
Write $[E]=\beta_E+\ddc\ell_E$ by the Poincar\'e--Lelong formula. Since $\ell_E$ is locally integrable and $T$ is smooth and $\ddc$-closed, distributional Stokes identifies the left side with $\int_X\beta_E\wedge T$. It is nonnegative for every such test, so $\kappa\geq0$. Also $D_p(T)>0$, which gives $\kappa<1$. Taking the supremum of $D_p(T)$ and using $\rho_\omega(p)=b^{-1/n}$ proves \eqref{eq:divisorial-coefficient-loss}.

If $c=b$ and $(T_j)$ is maximizing, the right side of the preceding measure bound tends to zero. Pairing with the fixed smooth form $\beta_E$ is continuous under current convergence, proving the cluster assertion. Equation~\eqref{eq:divisorial-coefficient-loss} forces $\kappa=0$, so \cref{cor:gauduchon-threshold} excludes bigness. Under the tail hypothesis, the near-maximizer statement in \cref{prop:integrable-threshold-adjoints} makes every cluster a maximizing current.
\end{proof}

At the threshold $c=b$, if $E\ne0$, the vanishing pairing defines a nonempty proper exposed face of the weak closure of smooth strictly positive normalized Gauduchon forms. Indeed, this closure is compact and convex; the pairing with $\beta_E$ is continuous and nonnegative there, and it is strictly positive on every smooth strictly positive member. The proposition places every cluster of a smooth maximizing sequence on this face. This assertion concerns that closure; it does not assert that every normalized Gauduchon current is a limit of smooth ones. Nor does the smooth pairing $\int_X\beta_E\wedge T$ define $[E]\wedge T$ for a singular $T$.

The hypotheses of the next two corollaries on $b$ and on the maximizing sequence hold whenever $p\in\mathcal D_\omega^0$, since this gives $b>0$ and uniform integrability over the whole of $\mathscr G_\omega$ (\cref{thm:integrable-current-attainment}).

\Needspace{12\baselineskip}
\begin{corollary}\label{cor:semipositive-divisor-exclusion}
Assume $b>0$ and that $D_p$ has a smooth maximizing sequence $(T_j)$ with uniformly integrable densities $p k_{T_j}$. Let $E$ be an effective real divisor whose class has a smooth semipositive representative $\beta_E$. If
\[
 \int_{\{\beta_E=0\}}f\,\mathrm{d}V=0,
\]
then no threshold subsolution $u$ for $f$ satisfies $\omega+\ddc u\geq[E]$.
\end{corollary}
\begin{proof}
Suppose such a subsolution exists. By \cref{prop:threshold-divisorial-constraint}, there is a maximizing current $T\in\mathscr G_\omega$ with $\int_X\beta_E\wedge T=0$. We use the matrix-density notation of \cref{sec:weak-gauduchon-duality}: $\sigma_T=\omega\wedge T$ is the trace measure of $T$, $P_T$ is its matrix density and $B$ is the coefficient matrix of $\beta_E$. Since both matrices are positive semidefinite, the vanishing pairing gives $\operatorname{tr}(BP_T)=0$ at $\sigma_T$-almost every point. Wherever $B\ne0$, this forces $P_T$ to be singular: a positive definite $P_T$ would have strictly positive trace against the nonzero positive semidefinite matrix $B$. The determinant measure $\Gamma_0(T)$ vanishes where $P_T$ is singular, so it is concentrated on $\{\beta_E=0\}$. By \cref{thm:gaud-determinant-subcritical}, $\Gamma_0(T)=k_T\,\mathrm{d}V$, and the hypothesis on $f$ gives $p k_T=0$ almost everywhere. This contradicts $D_p(T)=\rho_\omega(p)>0$.
\end{proof}


\begin{corollary}\label{cor:nef-divisor-threshold-volume}
Let $0\leq f\in L^1(\mathrm{d}V)$ be nonzero, with $f^{1/n}\in\mathcal D_\omega^0$. Suppose a subsolution at $b_\omega(f)$ satisfies $\omega+\ddc u\geq[E]$ for an effective real divisor $E$ whose Bott--Chern class is nef. Then
\[
 \lambda_\mu(\beta_E)=0
\]
for every smooth positive probability measure $\mu$, where $\beta_E$ is any smooth closed representative of that class.
\end{corollary}
\begin{proof}
By \cref{prop:threshold-divisorial-constraint}, there is a maximizing current $T\in\mathscr G_\omega$ with $\int_X\beta_E\wedge T=0$. If $\lambda_\mu(\beta_E)>0$, \cref{prop:nef-annihilator-determinant} gives $\Gamma_0(T)=0$. It follows that $D_{f^{1/n}}(T)=0$, contradicting its positive maximal value. Since nefness gives $\lambda_\mu(\beta_E)\geq0$, the conclusion follows for each $\mu$.
\end{proof}

This criterion uses nefness and the variational volume, without requiring a smooth semipositive representative. It complements \cref{cor:semipositive-divisor-exclusion}: the Hopf hyperplane class has zero variational volume, yet its nowhere vanishing semipositive representative excludes a component at the threshold. Neither argument identifies positive nef volume with bigness or classifies all possible divisorial components.


These are necessary conditions on divisorial singularities at the maximal global coefficient. They do not apply unchanged to a smaller coefficient or to a threshold defined within a prescribed singularity type. Determining which annihilating currents occur, and whether these conditions can be supplemented to characterize the possible divisorial components, remains a geometric realization problem. No identification with a null locus or a non-K\"ahler locus is asserted.


\section{Minimal multiplicities and the non-nef locus}
\label{sec:bc-minimal-multiplicities}

Return to the closed class $\alpha$ and assume it is pseudoeffective. Minimal potentials in the families $\alpha[-\varepsilon\omega]$ measure the singularities that persist as the allowed negative error tends to zero. Their Lelong numbers define the non-nef locus.

For $\varepsilon>0$ put
\[
 V_\varepsilon=V_{\theta+\varepsilon\omega},
 \qquad T_{\min,\varepsilon}=\theta+\ddc V_\varepsilon,
\]
using the normalized envelope of \cref{sec:obstacle-envelopes}. Pseudoeffectivity makes its competitor family nonempty. $V_\varepsilon$ has minimal singularities in this family: for any $u\in\PSH(X,\theta+\varepsilon\omega)$,
$u-\sup_Xu\leq V_\varepsilon$. Thus $T_{\min,\varepsilon}\in\alpha[-\varepsilon\omega]$.

\Needspace{10\baselineskip}
\begin{definition}\label{def:bc-minimal-multiplicity}
For $x\in X$, the \emph{minimal multiplicity} of $\alpha$ at $x$ is
\[
 \nu(\alpha,x)=\sup_{\varepsilon>0}\nu(T_{\min,\varepsilon},x).
\]
The \emph{non-nef locus} is
\[
 E_{\mathrm{nn}}(\alpha)=\{x\in X:\nu(\alpha,x)>0\}.
\]
\end{definition}

If $\varepsilon$ decreases, the admissible family shrinks, so its minimal multiplicities increase. Thus the supremum is the limit as $\varepsilon\downarrow0$. Comparison of two auxiliary metrics proves independence of $\omega$. Changing the representative changes minimal potentials only by a smooth function and bounded normalization errors, so it does not affect these numbers either.

The smooth regularization in \cref{subsec:psh-closed-regularization} turns vanishing Lelong numbers into an arbitrarily small curvature loss. This gives the converse in the following characterization.

\Needspace{10\baselineskip}
\begin{proposition}\label{prop:bc-non-nef}
For a pseudoeffective class $\alpha$, the set $E_{\mathrm{nn}}(\alpha)$ is a countable union of analytic subsets. Moreover,
\[
 \alpha\text{ is nef}\quad\Longleftrightarrow\quad
 E_{\mathrm{nn}}(\alpha)=\varnothing.
\]
\end{proposition}
\begin{proof}
By Siu's analyticity theorem from \cref{subsec:psh-lelong-levels} and monotonicity of the minimal multiplicities,
\[
 E_{\mathrm{nn}}(\alpha)
 =\bigcup_{k,\ell\geq1}
       \{x:\nu(T_{\min,1/\ell},x)\geq1/k\}.
\]
If $\alpha$ is nef, a smooth representative is available in every $\alpha[-\varepsilon\omega]$. A minimal potential is bounded below by its smooth potential up to a constant, so all its Lelong numbers vanish. Conversely, if the non-nef locus is empty, all $T_{\min,\varepsilon}$ have zero Lelong numbers. Apply the regularization consequence just stated to obtain a smooth representative bounded below by $-2\varepsilon\omega$. Letting $\varepsilon$ vary proves nefness.
\end{proof}

The parameter $\varepsilon$ is essential. Even for a nef class, a positive current with minimal singularities in $\alpha$ itself may have nonzero Lelong numbers; see the example discussed in \cref{sec:history-closed-loci}. \Cref{def:bc-minimal-multiplicity} does not replace the almost positive approximating families by the single family of positive currents.

\Needspace{13\baselineskip}
\section{The non-K\"ahler locus}
\label{sec:bc-non-kahler-locus}

For a big class, strict currents determine where the class admits a smooth positive representative. Analytic regularization and a finite-intersection argument produce one strict current whose polar set is exactly the obstruction.

For an almost positive closed $(1,1)$-current $T$ on $X$, write $E_+(T)=\{x:\nu(T,x)>0\}$.

\begin{definition}\label{def:bc-non-kahler-locus}
If $\alpha$ is big, let $T$ range over its K\"ahler-current representatives. Its \emph{non-K\"ahler locus} is
\[
 E_{\mathrm{nK}}(\alpha)=\bigcap_{T}E_+(T).
\]
For a pseudoeffective class that is not big, we use the convention $E_{\mathrm{nK}}(\alpha)=X$.
\end{definition}

\Needspace{10\baselineskip}
\begin{theorem}\label{thm:bc-non-kahler-locus}
Let $\alpha$ be big. There is a K\"ahler current with analytic singularities, smooth off its polar set, whose positive Lelong set is exactly $E_{\mathrm{nK}}(\alpha)$. In particular this locus is a proper analytic subset, and
\[
 E_{\mathrm{nn}}(\alpha)\subset E_{\mathrm{nK}}(\alpha),
 \qquad
 E_{\mathrm{nK}}(\alpha)=\varnothing
 \quad\Longleftrightarrow\quad\alpha\text{ is K\"ahler}.
\]
\end{theorem}
\begin{proof}
Analytic regularization reduces the defining intersection to K\"ahler currents with analytic singularities and smooth remainder away from their poles. Indeed, regularization of any strict current remains strict after a sufficiently small loss, and its positive Lelong set is contained in that of the original current.

Given two strict currents with potentials $u_1,u_2$, their maximum is a potential for a current dominating $\delta\omega$ for a common $\delta>0$, by \cref{lem:max-psh}. Its Lelong numbers are the minima of the two Lelong numbers. Regularizing this current gives a strict current with analytic singularities whose positive Lelong set is contained in the intersection of the previous two sets.

The descending chain condition for analytic subsets of a compact complex manifold now gives a smallest finite intersection of the analytic sets under consideration. This intersection equals the full intersection. The maximum-and-regularization construction realizes a subset of that finite intersection as the polar set of a single strict analytic current. Since the full intersection is contained in the polar set of every such current, equality follows.

For any positive current $T\in\alpha$ and any $\varepsilon>0$, minimality gives $\nu(T_{\min,\varepsilon},x)\leq\nu(T,x)$. Take the supremum in $\varepsilon$, then the intersection over strict $T$, to obtain the inclusion of loci. If the non-K\"ahler locus is empty, the realizing analytic current has no poles and is a smooth K\"ahler representative. The reverse implication follows by using a smooth K\"ahler form in the defining intersection.
\end{proof}

\Needspace{16\baselineskip}
\begin{proposition}\label{prop:bc-no-isolated-locus}
The non-K\"ahler locus of a big class has no isolated points.
\end{proposition}
\begin{proof}
Let $\theta+\ddc\Phi\geq\delta\omega$ be a current as in \cref{thm:bc-non-kahler-locus}. Suppose $x$ is an isolated point of its polar set. Choose holomorphic coordinates $z=(z_1,\ldots,z_n)$ on a ball centered at $x$, with no other pole. For a sufficiently large positive constant $A$ and a real constant $C$ to be chosen, the local function
$q(z)=A|z|^2-C$ satisfies $\theta+\ddc q\geq\delta'\omega$ on a smaller ball, for some $\delta'>0$.

Choose $C$ so that $q<\Phi$ on a closed annular collar of the boundary. This is possible because $\Phi$ is smooth on that collar. Since $\Phi\to-\infty$ at $x$, the maximum $\max(\Phi,q)$ agrees with $q$ near $x$ and with $\Phi$ on the collar. By \cref{lem:qpsh-gluing}, it glues to $\Phi$ outside the ball and gives a global K\"ahler current smooth near $x$. Its Lelong number at $x$ is zero, contradicting the definition of $E_{\mathrm{nK}}(\alpha)$.
\end{proof}

Closed bigness has a strong global consequence: the existence of a K\"ahler current is equivalent to membership in Fujiki's class $\mathcal C$, by the standard modification criterion recorded in \cref{sec:history-closed-loci}. Thus a compact complex manifold outside $\mathcal C$ has no big ordinary Bott--Chern class. The strictly positive nonclosed backgrounds of this book still exist on every compact complex manifold.

\subsection{Smooth equations on the ample locus}
\label{sec:ample-locus-smoothness}

Let $X$ be a connected compact complex manifold of dimension $n$, let $\alpha\in H^{1,1}_{\mathrm{BC}}(X,\mathbb R)$ be big, and let $\theta$ be a smooth closed representative. With the non-K\"ahler locus of \cref{def:bc-non-kahler-locus}, write
\[
 \operatorname{Amp}(\alpha)=X\setminus E_{\mathrm{nK}}(\alpha).
\]
By \cref{thm:bc-non-kahler-locus}, this is the locus on which a strict
current with analytic singularities can be smooth. We now assume that
$X$ is K\"ahler. The interior estimates of
\cref{sec:nonclosed-interior-smoothness} give the following application
to an arbitrary big class. The class is not required to be nef.

In the K\"ahler setting, $\operatorname{vol}(\alpha)$ denotes the classical volume of the big class. Write $\langle(\theta+\ddc u)^n\rangle$ for its non-pluripolar product, and $\mathcal E(X,\theta)$ for the $\theta$-plurisubharmonic potentials whose product has total mass $\operatorname{vol}(\alpha)$.

\begin{theorem}\label{thm:big-class-ample-smoothness}
Let $X$ be a compact K\"ahler manifold, let $\alpha$ be a big real
$(1,1)$-class, and let $\theta$ be a smooth closed representative.
Let $\mu$ be a smooth positive volume form with
$\mu(X)=\operatorname{vol}(\alpha)$. The normalized full-mass solution
$\varphi\in\mathcal E(X,\theta)$ of
\begin{equation}\label{eq:big-class-smooth-volume}
 \bigl\langle(\theta+\ddc\varphi)^n\bigr\rangle=\mu,
 \qquad \sup_X\varphi=0,
\end{equation}
is smooth on $\operatorname{Amp}(\alpha)$. The restriction of
$\theta+\ddc\varphi$ to this locus is a K\"ahler form.
\end{theorem}

Existence and uniqueness in this statement are the classical theorem
\cite[Theorem~A, equivalently Theorem~3.1]{BEGZ10}.
The nef case of the smoothness conclusion is
\cite[Theorem~C]{BEGZ10}; the following proof uses no nefness.

\begin{proof}
Fix a K\"ahler form $\omega$ and put $f=\mu/\omega^n$.
Compactness and smoothness give $f\in L^p(X,\omega^n)$ for every
finite $p>1$. Let $V_\theta=(\sup\{v\in\PSH(X,\theta):v\leq0\})^*$ be the normalized minimal envelope. The minimal-singularity estimate
\cite[Theorem~B, quantitatively Theorem~4.1]{BEGZ10} gives a constant $C>0$ such that
\begin{equation}\label{eq:ample-solution-minimal-bounds}
 V_\theta-C\leq\varphi\leq V_\theta,
 \qquad
 V_\theta=\left(\sup\{v\in\PSH(X,\theta):v\leq0\}\right)^*.
\end{equation}
The upper bound follows from the normalization.

Choose a strict analytic current $T$ of \cref{thm:bc-non-kahler-locus}. Since its Bott--Chern class is $\alpha$, it has a genuine potential $\rho$ with analytic singularities and smooth remainder. For some constant $\varepsilon>0$ it satisfies
\[
 T=\theta+\ddc\rho\geq\varepsilon\omega,
 \qquad \varepsilon>0,
 \qquad \{\rho=-\infty\}=E_{\mathrm{nK}}(\alpha),
\]
Thus $\rho$ is smooth on
$\Omega=\operatorname{Amp}(\alpha)$. Subtract a constant so that
$\rho\leq0$, hence $\rho\leq V_\theta$. By
\eqref{eq:ample-solution-minimal-bounds},
\[
 \rho-C\leq\varphi\leq0.
\]
Consequently $\varphi$ is locally bounded on $\Omega$.
On each relatively compact coordinate ball there, the truncations
which define the non-pluripolar product agree with $\varphi$ for all
sufficiently deep levels. Locality therefore identifies
\eqref{eq:big-class-smooth-volume} with the bounded-potential equation
\[
 (\theta+\ddc\varphi)^n=f\,\omega^n\quad\text{on }\Omega.
\]

Apply \cref{thm:big-minimal-holder}, the reconstructed local H\"older
result of \cite[Theorem~3.1]{Dang26}. Its hypotheses hold for this strict
barrier: the density is globally in $L^p$, and the solution has the
minimal bounds \eqref{eq:ample-solution-minimal-bounds}.
We obtain local H\"older continuity on all of $\Omega$.
\Cref{thm:nonclosed-interior-smoothness} now gives
$\varphi\in C^\infty(\Omega)$, since $f$ is smooth and strictly positive
there. The current $\theta+\ddc\varphi$ is a smooth semipositive form
on $\Omega$ with positive determinant $f\,\omega^n$, so every eigenvalue
is positive. It is closed because $\theta$ is closed. Hence it is a
K\"ahler form on $\Omega$.
\end{proof}

The proof of the H\"older input is given in
\cref{sec:big-reference-stability,sec:big-minimal-holder}, and the local
derivative estimates are proved in \cref{sec:nonclosed-interior-smoothness}.
The global $L^p$ hypothesis is automatic for the volume form in this
theorem. It cannot be replaced in that proof by smoothness only on the
ample locus. No extension of the metric across $E_{\mathrm{nK}}(\alpha)$
or completeness on its complement is asserted.


\section{Null loci and numerical restricted volumes}

Assume in this section that $X$ is K\"ahler, and choose $\omega$ K\"ahler. Closedness makes intersection numbers independent of representatives. For a nef class $\alpha$, their vanishing defines the classical null locus
\[
 \operatorname{Null}(\alpha)
 =\bigcup_{\substack{V\subset X\ \mathrm{irreducible},\ \dim V>0\\
                       \int_V\alpha^{\dim V}=0}}V.
\]
The intersection number is computed on a resolution of $V$ by pulling back a smooth closed representative; Stokes' theorem gives its independence of that representative. For a nef and big class, the equality
\[
 E_{\mathrm{nK}}(\alpha)=\operatorname{Null}(\alpha)
\]
is the transcendental Nakamaye theorem of Collins--Tosatti. Their theorem also applies on compact manifolds in class $\mathcal C$.

For a general pseudoeffective class on a compact K\"ahler manifold, the appropriate replacement is the numerical restricted volume in \eqref{eq:bc-numerical-restricted-volume}. For an irreducible compact analytic subvariety $V\subset X$ of dimension $k>0$, let $V_{\mathrm{reg}}$ be its regular locus. The numerical restricted volume is zero when $V\subset E_{\mathrm{nn}}(\alpha)$. Otherwise, for real $\varepsilon>0$, let $T$ range over closed representatives of $\alpha$ with $T\geq-\varepsilon\omega$, analytic singularities with smooth remainder, and polar set not containing $V$. The subscript $\mathrm{ac}$ denotes the absolutely continuous part on $V_{\mathrm{reg}}$. The volume is
\begin{equation}\label{eq:bc-numerical-restricted-volume}
 \langle\alpha^k\rangle_{X\mid V}
 =\lim_{\varepsilon\downarrow0}\ \sup_T
 \int_{V_{\mathrm{reg}}}
      \bigl((T+\varepsilon\omega)|_{V_{\mathrm{reg}}}\bigr)_{\mathrm{ac}}^k.
\end{equation}
The smooth remainder makes this measure the extension by zero of a smooth measure away from an analytic subset.

The numerical null locus is the union of the positive-dimensional irreducible subvarieties on which \eqref{eq:bc-numerical-restricted-volume} vanishes. The trace theory of \cite[Chapter 8]{SGPT} relates these restricted volumes to masses on $V$, which explains why singularity types and restriction of multiplier ideals are relevant to identifying this locus. Such a trace need not equal an ordinary pointwise restriction of a singular potential; see \cref{chap:trace-data}.

\section{Restricted volumes and Hermitian perturbations}

Return to an arbitrary compact complex $X$ and Hermitian metric $\omega$. Replacing the closed perturbation in \eqref{eq:bc-numerical-restricted-volume} by $\omega$ can give infinite volume even for the zero Bott--Chern class $\{0\}$. The unbounded smooth masses on the Hopf threefold already imply this failure in the ambient case $V=X$.

\Needspace{15\baselineskip}
\begin{proposition}\label{prop:bc-zero-class-infinite-error-volume}
Suppose there are smooth real functions $v_j$ on $X$, for positive integers $j$, such that
\[
 \omega+\ddc v_j>0,
 \qquad \int_X(\omega+\ddc v_j)^n\longrightarrow+\infty.
\]
For every $\varepsilon>0$,
\begin{equation}\label{eq:bc-zero-class-infinite-error-volume}
 \sup_{\substack{T\in\{0\}\ \mathrm{smooth,\ closed}\\
                          T\geq-\varepsilon\omega}}
       \int_X(T+\varepsilon\omega)^n=+\infty.
\end{equation}
This occurs for the Gauduchon metric on the Hopf threefold of \cref{chap:hopf}.
\end{proposition}
\begin{proof}
Fix $\varepsilon>0$ and take $T_j=\varepsilon\ddc v_j$. Each $T_j$ is smooth, closed and represents the zero Bott--Chern class. Moreover,
\[
 T_j+\varepsilon\omega
 =\varepsilon(\omega+\ddc v_j)>0,
\]
and its top mass is
$\varepsilon^n\int_X(\omega+\ddc v_j)^n\to+\infty$.
The Hopf specialization follows from \cref{thm:hopf-unbounded}.
\end{proof}

All currents in this proof have empty polar set, so imposing analytic singularities with smooth remainder does not eliminate the example. The supremum is already infinite for each fixed error size; letting $\varepsilon$ tend to zero afterwards does not repair it. In contrast, if $\omega$ is K\"ahler and $T=\ddc v$ is smooth, Stokes gives
\[
 \int_X(T+\varepsilon\omega)^n
 =\varepsilon^n\int_X\omega^n\longrightarrow0.
\]
\Needspace{7\baselineskip}
The zero Bott--Chern class is nef, has empty non-nef locus, and contains only the zero positive current: a positive representative would be $\ddc u\geq0$, so $u$ would be a global plurisubharmonic function on a compact manifold and hence constant. Thus \eqref{eq:bc-zero-class-infinite-error-volume} measures the nonlinear mass of the Hermitian error family, not positivity of the closed class.

A numerical restricted-volume theory in a Hermitian background must therefore control the approximating family or use another invariant. The qualitative non-Hermitian locus, defined through strict currents with a prescribed exterior derivative, does not use this supremum. Dang's qualitative null locus instead tests failure of bigness after restriction; the exact singularity condition and its equality with the non-Hermitian locus for nef and big backgrounds are recalled in \cref{sec:non-kahler-geometric-problems}.

\section{Finite restricted masses of closed classes}
\label{sec:finite-restricted-invariants}

The preceding infinite supremum comes from the nonclosed perturbation. For a closed class, a K\"ahler modification supplies a finite comparison mass. This gives an invariant on every subvariety, including those contained in the exceptional image of a chosen modification.

Assume in this section and the next that $X$ belongs to Fujiki's class $\mathcal C$. Every irreducible reduced subvariety $Y\subset X$ has a resolution $\pi:Z\to Y$ with $Z$ compact K\"ahler. Indeed, closed subspaces of a class-$\mathcal C$ space remain in that class, and the K\"ahler modification theorem applies to $Y$. These classical geometric inputs are recorded in \cref{sec:history-closed-loci}. In particular, taking the proper transform under one fixed modification of $X$ is not needed to obtain a model of $Y$.

Two such resolutions admit a common K\"ahler refinement: resolve the closure of their common graph in $Z_1\times Z_2$ by embedded resolution. The resulting smooth strict transform is K\"ahler, since the ambient resolution is a sequence of blowups of the compact K\"ahler product along smooth centers. The K\"ahler property of a blowup follows by adding a small positive multiple of a curvature form of $\mathcal O(-E)$ to the pulled-back metric; the curvature is positive in the vertical directions along the exceptional divisor $E$.

For a closed real $(1,1)$-class $b$ on a compact K\"ahler manifold $Z$ of dimension $k>0$, set
\begin{equation}\label{eq:closed-class-mass}
 \mathcal R_Z(b)=\sup_{\substack{S\in b\\S\geq0}}
                  \int_Z\langle S^k\rangle,
 \qquad \sup\varnothing=0.
\end{equation}
For a class $a\in H^{1,1}_{\mathrm{BC}}(X,\mathbb R)$ and an irreducible reduced $k$-dimensional $Y\subset X$, write $f=j\circ\pi$, where $j:Y\hookrightarrow X$, and define
\begin{equation}\label{eq:finite-restricted-invariant}
 \mathcal R(a;Y)=\mathcal R_Z(f^*a).
\end{equation}

\begin{proposition}\label{prop:finite-restricted-invariance}
The quantities in \eqref{eq:closed-class-mass} and \eqref{eq:finite-restricted-invariant} are finite. The latter is independent of the K\"ahler resolution. If $\eta\in b$ is smooth, $\kappa$ is K\"ahler and $C>0$ satisfies $C\kappa-\eta>0$, then
\begin{equation}\label{eq:closed-class-mass-bound}
 \mathcal R_Z(b+t\{\kappa\})
       \leq(C+t)^k\int_Z\kappa^k \qquad(t\geq0).
\end{equation}
For a modification $\mu:X'\to X$ and the proper transform $Y'$ of a subvariety meeting the isomorphism locus,
\[
 \mathcal R(\mu^*a;Y')=\mathcal R(a;Y).
\]
\end{proposition}
\begin{proof}
First consider a K\"ahler form $\lambda$ on $Z$. For $q\in\PSH(Z,\lambda)$ normalized by $q\leq0$, the truncations $q_j=\max(q,-j)$ have mass $\int_Z\lambda^k$ by \cref{thm:mass-invariance-characterization}, or Stokes' theorem when $k=1$. Restricting their measures to $\{q>-j\}$ and passing to the increasing limit gives
\[
 \int_Z\langle(\lambda+\ddc q)^k\rangle
       \leq\int_Z\lambda^k.
\]
The smooth current $\lambda$ attains equality.

Put $A=C\kappa-\eta$. For a positive current $S$ in $b+t\{\kappa\}$, the current $S+A$ is positive and belongs to $(C+t)\{\kappa\}$. Positivity and multilinearity of local non-pluripolar products give $\langle S^k\rangle\leq\langle(S+A)^k\rangle$. The K\"ahler bound just proved yields \eqref{eq:closed-class-mass-bound}. The definition uses the class itself, so changing its smooth representative or the comparison form does not change the supremum.

Let $\rho:W\to Z$ be a modification between compact K\"ahler manifolds. By \cref{prop:psh-descent-modification}, every positive current in $\rho^*b$ is the pullback of a unique positive current in $b$. Their non-pluripolar measures correspond on the isomorphism locus, and both give zero mass to the exceptional analytic sets. Hence their total masses agree and
\[
 \mathcal R_W(\rho^*b)=\mathcal R_Z(b).
\]
Apply this identity on a common K\"ahler refinement to obtain independence of resolution. The same argument compares resolutions of $Y$ and $Y'$, since their maps to $X$ agree on a dense open set and therefore everywhere on the common refinement.
\end{proof}

Positivity of this mass has a classical interpretation on a K\"ahler model. We use two facts from K\"ahler pluripotential theory: a class containing a positive current with positive non-pluripolar top mass is big, and a current with minimal singularities in a big class attains its positive-product volume. For a nef class this volume is its top intersection. Their precise sources are recorded in \cref{sec:history-closed-loci}. The strict current needed for a nef class can also be obtained by the following normalized diagonal construction.

\begin{proposition}\label{prop:restricted-mass-bigness}
For a closed class $b$ on a compact K\"ahler $k$-fold $Z$,
\[
 \mathcal R_Z(b)>0\quad\Longleftrightarrow\quad b\text{ is big}.
\]
If $b$ is nef, then
\begin{equation}\label{eq:nef-restricted-mass}
 \mathcal R_Z(b)=\int_Zb^k.
\end{equation}
For nef $b$ with $I=\int_Zb^k>0$, fix a K\"ahler form $\kappa$, let $p_1,p_2:Z\times Z\to Z$ be the projections and $\Delta\subset Z\times Z$ the diagonal, and put
\[
 M=k\int_Zb^{k-1}\wedge\{\kappa\}.
\]
There is a closed positive $(k,k)$-current $\Theta$ in $(p_1^*b+p_2^*b)^k$ with $\Theta\geq d[\Delta]$, $d>0$, such that
\begin{equation}\label{eq:normalized-diagonal-current}
 K_b=M^{-1}(p_1)_*(\Theta\wedge p_2^*\kappa)\in b,
 \qquad K_b\geq(d/M)\kappa.
\end{equation}
\end{proposition}
\begin{proof}
If $\mathcal R_Z(b)>0$, some positive representative has positive top mass, so the positive-mass criterion makes $b$ big. Conversely, if $S\in b$ satisfies $S\geq\delta\kappa$, then $S-\delta\kappa$ is positive and closed. Expansion of the non-pluripolar product gives $\langle S^k\rangle\geq\delta^k\kappa^k$.

Suppose $b$ is nef. For $t>0$, the class $b+t\{\kappa\}$ is K\"ahler. Adding $t\kappa$ to an arbitrary positive representative of $b$ and using the K\"ahler mass bound yields
\[
 \mathcal R_Z(b)\leq\int_Z(b+t\{\kappa\})^k.
\]
Letting $t$ decrease to zero gives the upper bound in \eqref{eq:nef-restricted-mass}. In particular, equality holds if $I=0$.

For $I>0$, choose $L$ so that $L\{\kappa\}-b$ is K\"ahler. Nonnegativity of intersections of nef classes gives
\[
 0\leq\int_Zb^{k-1}\wedge(L\{\kappa\}-b)
      =LM/k-I,
\]
so $M>0$. The class $B=p_1^*b+p_2^*b$ is nef and satisfies
\[
 \int_{Z\times Z}B^{2k}={2k\choose k}I^2>0.
\]
The classical mass-concentration theorem along the diagonal gives $\Theta\in B^k$ with $\Theta\geq d[\Delta]$ for some $d>0$. In pushing forward $B^k\wedge p_2^*\{\kappa\}$, only the term with one factor from the first copy of $Z$ survives. Thus
\[
 (p_1)_*(B^k\wedge p_2^*\{\kappa\})=Mb,
 \qquad
 (p_1)_*([\Delta]\wedge p_2^*\kappa)=\kappa.
\]
This proves \eqref{eq:normalized-diagonal-current}. The K\"ahler $\ddc$-lemma for currents identifies the computed cohomology class with the Bott--Chern class $b$. Therefore $b$ is big, and a current with minimal singularities has mass $\int_Zb^k$. This gives the remaining lower bound.
\end{proof}

The criterion descends to a smooth subvariety $Y$. Write $a_Y=j^*a$. If $\pi:Z\to Y$ is a K\"ahler modification and $K\in\pi^*a_Y$ satisfies $K\geq\delta\kappa$, choose a Hermitian metric $h$ on $Y$ and $C>0$ with $\pi^*h\leq C\kappa$. Potential descent gives a positive representative $\pi_*K$ of $a_Y$, and
\begin{equation}\label{eq:strict-current-pushforward}
 \pi_*K\geq(\delta/C)\pi_*\pi^*h=(\delta/C)h.
\end{equation}
Conversely, a strict positive representative of $a_Y$ pulls back to a positive current whose non-pluripolar mass is positive on a coordinate ball in the isomorphism locus. On that ball it dominates a positive multiple of the closed form $\kappa$, so the local product expansion applies. Consequently, $\mathcal R(a;Y)>0$ is equivalent to bigness of the restricted class $a_Y$. This argument uses the Hermitian metric only for the current inequalities; the product comparison takes place with a closed local form.

\begin{proposition}\label{prop:restricted-rank-drop}
Let $W$ be a connected compact K\"ahler manifold of dimension $k>0$, let $g:W\to X$ be holomorphic, and suppose $\dim g(W)<k$. Then for every closed class $a$ on $X$,
\[
 \mathcal R_W(g^*a)=0.
\]
\end{proposition}
\begin{proof}
If the image is a point, $g^*a=0$ and \eqref{eq:nef-restricted-mass} gives the assertion. Otherwise, choose a K\"ahler resolution $h:Q\to g(W)$, and resolve the graph of the meromorphic lift of $g$. This gives a K\"ahler modification $\rho:\widehat W\to W$ and a holomorphic map $F:\widehat W\to Q$ with $g\circ\rho=h\circ F$, where $h$ is also viewed as a map into $X$.

Choose a smooth representative $\eta\in a$, a K\"ahler form $\kappa_Q$ on $Q$, and $C>0$ such that $C\kappa_Q-h^*\eta\geq0$. The forms
\[
 P=CF^*\kappa_Q,
 \qquad A=P-\rho^*g^*\eta
\]
are closed and semipositive. Since $\dim Q<k$, $P^k=0$. For any positive current $S\in g^*a$, the current $\rho^*S+A$ is positive in the nef class $\{P\}$. Hence
\[
 \int_{\widehat W}\langle(\rho^*S)^k\rangle
 \leq\int_{\widehat W}\langle(\rho^*S+A)^k\rangle
 \leq\mathcal R_{\widehat W}(\{P\})
 =\int_{\widehat W}P^k=0.
\]
Modification invariance proves the result. If the class has no positive representative, its invariant is zero by definition.
\end{proof}

\begin{corollary}\label{cor:finite-restricted-null-union}
Let $a$ be nef on $X$. For every positive-dimensional irreducible reduced $Y\subset X$,
\[
 \mathcal R(a;Y)=\int_Ya^{\dim Y}.
\]
If $\int_Xa^n>0$, then
\begin{equation}\label{eq:finite-restricted-null-union}
 \bigcup_{\substack{Y\subset X\ \mathrm{irreducible},\ \dim Y>0\\
                         \mathcal R(a;Y)=0}}Y
       =\operatorname{Null}(a)=E_{\mathrm{nK}}(a).
\end{equation}
\end{corollary}
\begin{proof}
Pullback preserves nefness: pull back smooth representatives with arbitrarily small negative error and bound the pulled-back Hermitian metric by a K\"ahler metric on the resolution. Now apply \eqref{eq:nef-restricted-mass}. The intersection number on $Y$ is defined by this same pullback. The first equality in \eqref{eq:finite-restricted-null-union} follows from the definition of the null locus; the second is the classical nef-and-big locus theorem recalled earlier in this chapter.
\end{proof}

When $X$ is K\"ahler and $a$ is nef, \eqref{eq:bc-numerical-restricted-volume} is also $\int_Ya^{\dim Y}$, so the two restricted invariants agree. For a general class, \eqref{eq:finite-restricted-invariant} takes the supremum over all positive currents on a resolution of $Y$; it does not require those currents to extend to $X$. The nef comparison therefore does not identify the two constructions outside that regime.


\section{Divisorial traces with bounded remainder}
\label{sec:resolved-trace-masses}

The invariant $\mathcal R(a;Y)$ is a supremum over all positive currents in a class. The trace of a given potential on $Y$, on the other hand, carries a specified singularity, and its non-pluripolar mass can be smaller. Suppose that on a resolution the trace becomes a divisorial logarithmic term plus a bounded remainder. The non-pluripolar product then ignores the divisor and sees only the residual current, which has bounded potentials. The class of this residual current is nef, and its top intersection number computes the mass.

We keep the assumption that $X$ belongs to class $\mathcal C$. Let $Y\subset X$ be an irreducible reduced subvariety of dimension $k>0$, let $\alpha_0$ be a smooth closed real $(1,1)$-form on $X$, and let $u\in\PSH(X,\alpha_0)$ be such that its pointwise restriction $v=u|_Y$ is not identically $-\infty$. On a K\"ahler resolution $\pi\colon Z\to Y$, with $f=j\circ\pi\colon Z\to X$ as in \cref{sec:finite-restricted-invariants}, put $\alpha_Z=f^*\alpha_0$ and $v_\pi=u\circ f$. When $Y$ is smooth, one may instead start from a smooth closed real $(1,1)$-form $\alpha_Y$ on $Y$ and $v\in\PSH(Y,\alpha_Y)$, and put $\alpha_Z=\pi^*\alpha_Y$ and $v_\pi=v\circ\pi$. In both cases $v_\pi$ is the composition of a given potential with a map; the intrinsic multiplier-ideal trace of \cref{chap:trace-data} is a different construction.

Suppose that on one such resolution
\begin{equation}\label{eq:resolved-trace-decomposition}
 v_\pi=\ell+w,
 \qquad \ell=\sum_{j=1}^N c_j\log|s_j|_{h_j}^2,
 \qquad w\text{ bounded},
\end{equation}
where $c_j>0$, the $D_j$ are prime divisors on $Z$, $s_j$ is the canonical section of $\mathcal O(D_j)$, and $h_j$ is a smooth Hermitian metric on $\mathcal O(D_j)$; the empty sum $\ell=0$ is allowed. Write $\Gamma=\sum_jc_jD_j$ and $\gamma_j=c_1(\mathcal O(D_j),h_j)$ for the curvature forms, and define the \emph{residual form} and the \emph{residual class}
\begin{equation}\label{eq:resolved-trace-residual-class}
 \eta_v=\alpha_Z-\sum_jc_j\gamma_j,
 \qquad b_v=\{\eta_v\}=\{\alpha_Z\}-\{[\Gamma]\}.
\end{equation}
For singular $Y$, all assertions below are made on its resolutions.

\begin{lemma}\label{lem:bounded-residual-mass-polynomial}\index{Residual class!mass on a resolution}
The current $R_v=\eta_v+\ddc w$ is positive and has bounded local potentials, and its class $b_v$ is nef. Moreover,
\begin{equation}\label{eq:resolved-trace-mass}
 \int_Z\langle(\alpha_Z+\ddc v_\pi)^k\rangle
       =\int_ZR_v^k
       =\int_Zb_v^k
       =\mathcal R_Z(b_v).
\end{equation}
\Needspace{7\baselineskip}
If $v_1,\ldots,v_k$ satisfy the same conditions on $Z$, possibly for different closed backgrounds, then
\begin{equation}\label{eq:resolved-trace-mixed-mass}
 \int_ZR_{v_1}\wedge\cdots\wedge R_{v_k}
       =\int_Zb_{v_1}\wedge\cdots\wedge b_{v_k}.
\end{equation}
\end{lemma}
\begin{proof}
Off the support $|\Gamma|$, we have $\ddc\ell=-\sum_jc_j\gamma_j$, so the current $\alpha_Z+\ddc v_\pi$ equals $\eta_v+\ddc w$ there and is positive. After the addition of a smooth local potential for $\eta_v$, the bounded function $w$ is plurisubharmonic off $|\Gamma|$. By \cref{lem:qpsh-removable-analytic}, it extends across the divisor as a bounded plurisubharmonic function that agrees almost everywhere with the original one; we use this extension to represent $w$. Thus $R_v\geq0$ on $Z$, and the Poincar\'e--Lelong formula gives
\[
 \alpha_Z+\ddc v_\pi=[\Gamma]+R_v.
\]
Since $w$ is bounded, $b_v$ is nef by \cref{lem:bounded-closed-potential-nef}. The non-pluripolar measure on the left of \eqref{eq:resolved-trace-mass} agrees with $R_v^k$ off the divisor, and neither measure charges the divisor. This gives the first equality.

Fix a K\"ahler form $\kappa$ on $Z$. For all sufficiently large real $t$, the smooth closed form $\eta_v+t\kappa$ is K\"ahler, and $w$ is a bounded potential for it. Therefore
\[
 \int_Z(R_v+t\kappa)^k=\int_Z(\eta_v+t\kappa)^k
\]
by \cref{thm:mass-invariance-characterization}, or by Stokes' theorem when $k=1$. Both sides are polynomials in $t$, so their constant coefficients agree; this gives $\int_ZR_v^k=\int_Zb_v^k$. The last equality in \eqref{eq:resolved-trace-mass} is \eqref{eq:nef-restricted-mass}. Applying the same polynomial argument to positive linear combinations of the $R_{v_i}$ and comparing coefficients gives \eqref{eq:resolved-trace-mixed-mass}.
\end{proof}

We denote the common value in \eqref{eq:resolved-trace-mass} by $\mathfrak r(v;Y)$ and call it the \emph{trace mass} of $v$ on $Y$. When $Y$ is smooth, it is also the non-pluripolar mass of the trace current on $Y$ itself:
\[
 \mathfrak r(v;Y)=\int_Y\langle(\alpha_Y+\ddc v)^k\rangle,
\]
where $\alpha_Y=\alpha_0|_Y$ when $v$ is a restriction. Indeed, the currents on $Z$ and on $Y$ correspond on the locus where $\pi$ is an isomorphism, and their non-pluripolar measures give no mass to its analytic complement. A change of variables therefore identifies the total masses.

\begin{theorem}\label{thm:resolved-trace-invariance}\index{Restricted volume!resolution independence}
The number $\mathfrak r(v;Y)$ does not depend on the divisor metrics or on the K\"ahler resolution carrying \eqref{eq:resolved-trace-decomposition}. If two potentials in the same closed background have the same ordinary singularity type and each has a decomposition of the form \eqref{eq:resolved-trace-decomposition}, then their residual classes agree on a common refinement, and in particular their trace masses are equal. Residual classes add when the traces and their backgrounds are added, and their mixed masses are given by \eqref{eq:resolved-trace-mixed-mass}.
\end{theorem}
\begin{proof}
Changing a divisor metric changes $\ell$ by a smooth function, $w$ by its negative, and $\eta_v$ by the corresponding $\ddc$-term. Thus $R_v$ is unchanged. The same cancellation occurs when a smooth change of the background is accompanied by the opposite change of the potential.

Let $\rho\colon W\to Z$ be a modification with $W$ K\"ahler. Pulling back the decomposition, with the total divisor $\rho^*\Gamma$ including all exceptional coefficients, gives a decomposition on $W$ with bounded remainder $w\circ\rho$, residual current $\rho^*R_v$ and nef residual class $\rho^*b_v$. Its mass is
\[
 \int_W(\rho^*b_v)^k=\int_Zb_v^k.
\]
For two decompositions on different resolutions, pass to a common K\"ahler refinement and pull both back in this way. At a generic smooth point of a prime component of the union of the two divisors, the difference of the logarithmic terms is $(a_1-a_2)\log|t|^2+O(1)$, where $t$ is a local defining coordinate of the component and $a_1,a_2$ are its two coefficients. The two potentials have bounded difference, and the resulting pair of inequalities forces $a_1=a_2$. Hence the total divisors coincide, and so do the residual classes in the common background. Finally, adding two decompositions adds both the divisors and the bounded remainders. The residual class of the sum is therefore the sum of the nef residual classes, and \eqref{eq:resolved-trace-mixed-mass} computes its mixed terms.
\end{proof}

\begin{proposition}\label{prop:resolved-trace-strict-current}
Suppose that $\mathfrak r(v;Y)>0$. Then the class $\{\alpha_Z\}$ contains a K\"ahler current $\alpha_Z+\ddc p$ whose potential satisfies $p\leq v_\pi+A$ for some constant $A$. If $b_v$ has a smooth positive representative on a resolution carrying \eqref{eq:resolved-trace-decomposition}, then on that resolution $p$ can be chosen with $p\sim v_\pi$, that is, with $p-v_\pi$ bounded. When $Y$ is smooth, the K\"ahler current and the corresponding inequalities between potentials descend to $Y$.
\end{proposition}
\begin{proof}
Positive trace mass means $\int_Zb_v^k>0$. Since $b_v$ is nef, \eqref{eq:normalized-diagonal-current} gives a current $K=\eta_v+\ddc a\geq\delta\kappa$ in $b_v$, where $\kappa$ is a K\"ahler form and $\delta>0$. The current
\[
 [\Gamma]+K=\alpha_Z+\ddc(\ell+a)
\]
is then a K\"ahler current. Since $a$ is bounded above on $Z$ and $w$ is bounded below, its potential $p=\ell+a$ satisfies $p\leq\ell+w+A=v_\pi+A$ for a constant $A$. If $b_v$ has a smooth positive representative, $a$ can be chosen smooth, and then $p-v_\pi=a-w$ is bounded.

For smooth $Y$, take $\alpha_Y=\alpha_0|_Y$ when $v$ is a restriction. By \cref{prop:psh-descent-modification}, $p=q\circ\pi$ for some $q\in\PSH(Y,\alpha_Y)$. The current $\alpha_Y+\ddc q$ is strict by \eqref{eq:strict-current-pushforward}, and since $\pi$ is surjective, each inequality between potentials descends. The hypothesis of a smooth positive representative refers to the specified resolution: the pullback of a K\"ahler class under a further blowup need not be K\"ahler.
\end{proof}

The trace mass can vanish even when the restricted class is big, because the non-pluripolar product ignores the divisorial part of the trace.

\begin{example}\label{ex:divisorial-trace-zero-mass}
Let $k\geq1$, let $Y\simeq\mathbb P^k$ be a hyperplane in $X=\mathbb P^{k+1}$, and let another hyperplane $H_X\subset X$ cut out $H=H_X\cap Y$. Let $\omega_{\mathrm{FS}}$ be the Fubini--Study form, normalized by $\int_Y\omega_{\mathrm{FS}}^k=1$, and let $s_{H_X}$ be a section of $\mathcal O(1)$ defining $H_X$, measured with the Fubini--Study metric. For $c>0$ and $d\geq0$, take $\alpha_0=(c+d)\omega_{\mathrm{FS}}$ on $X$ and $u=c\log|s_{H_X}|^2$. The restriction $v=u|_Y$ is not identically $-\infty$, and
\[
 \alpha_0|_Y+\ddc v=c[H]+d\omega_{\mathrm{FS}}|_Y,
 \qquad \mathfrak r(v;Y)=d^k,
 \qquad \mathcal R(\{\alpha_0\};Y)=(c+d)^k.
\]
Thus for $d=0$ the trace mass is zero, although the restricted invariant is positive. The divisor $c[H]$ is a nonzero part of the trace current, but the non-pluripolar top product discards it.
\end{example}

The next example shows that the inequality $p\leq v_\pi+A$ in \cref{prop:resolved-trace-strict-current} cannot in general be improved to $p\sim v_\pi$, even for a bounded trace.

\Needspace{14\baselineskip}
\begin{example}\label{ex:positive-trace-bounded-type-obstruction}
Let $\mu\colon Y\to\mathbb P^2$ be the blowup of a point, with exceptional curve $E$, let $\omega_{\mathrm{FS}}$ be the Fubini--Study form on $\mathbb P^2$ with total volume one, and put $\alpha_Y=\mu^*\omega_{\mathrm{FS}}$ and $v=0$. The decomposition \eqref{eq:resolved-trace-decomposition} holds on $Z=Y$ with no divisor, the residual class $\{\alpha_Y\}$ is nef and big, and
\[
 \mathfrak r(0;Y)=\mathcal R(\{\alpha_Y\};Y)=1.
\]
Nevertheless, there is no bounded $q\in\PSH(Y,\alpha_Y)$ for which $\alpha_Y+\ddc q\geq\delta\kappa$ with $\delta>0$ and $\kappa$ K\"ahler. Indeed, since $q$ is bounded, its restriction to $E$ is defined, and the disk criterion of \cref{lem:smooth-disk-criterion} gives
\[
 \ddc(q|_E)\geq\delta\kappa|_E,
\]
because $\alpha_Y|_E=0$. Integration over the compact curve $E$ contradicts Stokes' theorem. Thus positive trace mass gives the one-sided comparison of \cref{prop:resolved-trace-strict-current}, but a K\"ahler current in the class may need additional singularities.
\end{example}

A subvariety contained in the null locus can nevertheless have positive restricted invariant; the union in \eqref{eq:finite-restricted-null-union} is needed to recover the locus.

\begin{example}\label{ex:positive-restriction-inside-null-locus}
Let $\mu\colon X\to\mathbb P^3$ be the blowup along a line $L$, and set $a=\mu^*\{\omega_{\mathrm{FS}}\}$, where $\omega_{\mathrm{FS}}$ has total volume one; the class $a$ is nef and $\int_Xa^3=1$. Since $N_{L/\mathbb P^3}\simeq\mathcal O(1)\oplus\mathcal O(1)$, the exceptional divisor is $E=\mathbb P(N_{L/\mathbb P^3})\simeq\mathbb P^1\times\mathbb P^1$, and on $E$ the map $\mu$ is the projection to $L$. Hence $\int_Ea^2=0$, whereas a section $Y$ of $E\to L$ satisfies $\int_Ya=1$. Every irreducible subvariety not contained in $E$ has positive intersection number, because $\mu^*\omega_{\mathrm{FS}}$ is positive on its nonempty open part outside $E$. By \cref{cor:finite-restricted-null-union},
\[
 E_{\mathrm{nK}}(a)=E,
 \qquad \mathcal R(a;E)=0,
 \qquad \mathcal R(a;Y)=1.
\]
The union of the subvarieties with zero invariant is the whole null locus $E$, although the curve $Y\subset E$ has positive invariant. Positivity along one subvariety does not place it outside the null locus.
\end{example}


The elliptic example below lies outside class $\mathcal C$: every pseudoeffective closed class is nef and none is big, while the Hermitian equation has smooth positive solutions. Its explicit cones and degeneration orders allow these notions of positivity to be compared on one manifold.

\section{Closed currents on the elliptic threefold}
\label{sec:elliptic-closed-current-descent}

The lower-mass example also permits an explicit description of its positive closed currents. This gives a useful comparison between the Hermitian equation and the theory of closed Bott--Chern classes on the same manifold.

Let $X$ be the threefold of \cref{sec:vanishing-lower-mass}, and let
\[
 \pi:X\longrightarrow B=\mathbf P^1\times\mathbf P^1
\]
be its elliptic fibration. Write $\sigma_1,\sigma_2$ for the integral-one Fubini--Study forms from the two factors, so $\alpha=\pi^*\sigma_1$ and $\beta=\pi^*\sigma_2$. The product coordinates used in that construction show that $\pi$ is locally a holomorphic product with the connected elliptic curve $E$. We retain its balanced metric $\omega=\alpha+\beta+D$, where $D=\mathrm{d}t\wedge\dc t$. It satisfies
\[
 \alpha\wedge\beta=\tfrac12\ddc\omega.
\]
The argument below uses this identity, not the value of the lower Monge--Amp\`ere mass.

\subsection{Descent of positive closed currents}

\begin{lemma}\label{lem:elliptic-positive-current-descent}
Every positive closed real $(1,1)$-current $T$ on $X$ has the form
\[
 T=\pi^*S
\]
for a unique positive closed $(1,1)$-current $S$ on $B$. Pullback here is the ordinary pullback by a smooth holomorphic submersion.
\end{lemma}
\begin{proof}
The measure $T\wedge\alpha\wedge\beta$ is positive. Current Stokes gives
\[
 \int_XT\wedge\alpha\wedge\beta
   =\tfrac12\int_X\omega\wedge\ddc T=0.
\]
Thus this measure is zero. All products in this identity pair a current with smooth forms.
\Needspace{10\baselineskip}
In local product coordinates $(z_1,z_2,\zeta)$, write the coefficient measures of $T$ as densities $p_{j\bar k}$ of a positive Hermitian matrix relative to their trace measure. The smooth form $\alpha\wedge\beta$ is a strictly positive volume form in the two base directions. Its wedge with $T$ is a positive smooth factor times the vertical diagonal coefficient. That coefficient therefore vanishes. Positivity of the matrix and
$|p_{j\bar\zeta}|^2\leq p_{j\bar j}p_{\zeta\bar\zeta}$ then make its entire vertical row and column zero. Consequently $T$ has only terms of the form
\[
 \frac{\mathrm{i}}{2\pi}\,t_{j\bar k}(z,\zeta)
       \mathrm{d}z_j\wedge\mathrm{d}\bar z_k,
       \qquad 1\leq j,k\leq2,
\]
where the coefficients are distributions of order zero.

The equations $\partial T=\bar\partial T=0$ imply
\[
 \partial_\zeta t_{j\bar k}
       =\partial_{\bar\zeta}t_{j\bar k}=0
\]
in the sense of distributions: these derivatives are exactly the coefficients of the terms containing one vertical differential. A distribution on a product whose derivatives in both real fiber directions vanish is pulled back from the base. For completeness, convolve it on a smaller product polydisk. The convolutions have both fiber derivatives zero and are independent of the fiber variable. Testing them against a fixed fiber test function of integral one gives distributions on the base; passage to the limit identifies the original distribution with their pullback. Independence on overlapping fiber disks and connectedness of the fiber extend this conclusion across the whole local product $U\times E$.

We therefore have $T|_{\pi^{-1}(U)}=\pi^*S_U$. Positivity and closedness of $S_U$ follow by testing the coefficient distributions against a nonnegative fiber test function of integral one and using the horizontal equations. Pullback on distributions is injective in a product, by the same test. Hence the $S_U$ agree on overlaps and define the required positive closed current $S$. Injectivity also proves uniqueness.
\end{proof}

The Gauduchon forms in \cref{lem:lower-mass-unbounded-root} have bidegree $(2,2)$ and satisfy $\ddc\Gamma_j=0$. Their bounded degree and divergent positive-root volume are compatible with this description of the closed $(1,1)$-currents. The two statements concern different cones of currents.

\begin{proposition}\label{prop:elliptic-bott-chern-cones}
Pullback
\[
 \pi^*:H^{1,1}_{\mathrm{BC}}(B,\mathbb R)
           \longrightarrow H^{1,1}_{\mathrm{BC}}(X,\mathbb R)
\]
is injective. The pseudoeffective and nef cones on $X$ coincide and are
\[
 \operatorname{Psef}(X)=\operatorname{Nef}(X)
 =\{a\{\alpha\}_{\mathrm{BC}}+b\{\beta\}_{\mathrm{BC}}:
                                           a,b\geq0\}.
\]
Its trace-one base for the Gauduchon metric $\omega$ is the line segment $a+b=(8\pi)^{-1}$. Every class in this cone has a smooth semipositive representative, and none is big. The real form $\beta-\alpha$ is de Rham exact but represents a nonzero Bott--Chern class.
\end{proposition}
\begin{proof}
Suppose a smooth closed real form $\tau$ on $B$ satisfies $\pi^*\tau=\ddc u$ with $u$ smooth on $X$. Restriction to each compact connected elliptic fiber makes $u$ harmonic, hence constant there. Local holomorphic sections show that $u=v\circ\pi$ for a smooth function $v$ on $B$. Thus $\pi^*(\tau-\ddc v)=0$, which implies $\tau=\ddc v$ because $\pi$ is a submersion. This proves injectivity.

The standard cohomology calculation for $\mathbf P^1\times\mathbf P^1$, together with the K\"ahler $\partial\bar\partial$ lemma, gives
\[
 H^{1,1}_{\mathrm{BC}}(B,\mathbb R)
       =\mathbb R\{\sigma_1\}\oplus\mathbb R\{\sigma_2\},
 \qquad \int_B\sigma_1\wedge\sigma_2=1,
 \qquad \sigma_1^2=\sigma_2^2=0.
\]
These are the usual projective-space and product computations recalled in the closed theory. If $S$ is positive and closed, its two class coefficients are
\[
 a=\int_BS\wedge\sigma_2\geq0,
 \qquad b=\int_BS\wedge\sigma_1\geq0.
\]
The current description of Bott--Chern cohomology and the preceding lemma show that every pseudoeffective class on $X$ is therefore in the displayed cone. Conversely every such class has the smooth semipositive representative $a\alpha+b\beta$, and is nef. A nef class is pseudoeffective by the compactness argument following \cref{prop:psef-cone-compact-base}. This proves the equality of cones. It does not require a calculation of the whole Bott--Chern group of $X$.

Since $\omega^2=2\alpha\wedge\beta+2(\alpha+\beta)\wedge D$ and $\int_E D=4\pi$, fiber integration gives
\[
 \int_X(a\alpha+b\beta)\wedge\omega^2=8\pi(a+b).
\]
This proves the assertion about the trace-one base. No closed K\"ahler current exists on $X$, by \cref{prop:lower-mass-no-pluriclosed}, so none of these classes is big. Finally, the globally defined one-form $\dc t$ satisfies
\[
 \mathrm{d}(\dc t)=\beta-\alpha.
\]
The class of $\sigma_2-\sigma_1$ on $B$ is nonzero, and injectivity makes its pullback nonzero in Bott--Chern cohomology. Thus de Rham exactness does not imply $\ddc$-exactness on this example.
\end{proof}

In particular, every pseudoeffective class on this threefold has empty non-nef locus, whereas its non-K\"ahler locus is all of $X$ under the convention for non-big classes. The example does not contradict a nef-and-big theorem: no closed nef class here has positive top self-intersection. Indeed $(a\alpha+b\beta)^3=0$.

\Needspace{19\baselineskip}
\subsection{Potentials and envelopes from the base}

\begin{proposition}\label{prop:elliptic-potential-descent}
Let $a,b\geq0$, put $\theta_B=a\sigma_1+b\sigma_2$ and $\theta=\pi^*\theta_B$. Pullback is a bijection
\[
 \PSH(B,\theta_B)\longrightarrow\PSH(X,\theta),\qquad v\longmapsto v\circ\pi.
\]
It preserves ordinary singularity order and equivalence, and ideal singularity order and equivalence. For every $v\in\PSH(B,\theta_B)$,
\[
 P_\theta[v\circ\pi]=P_{\theta_B}[v]\circ\pi,
 \qquad Q_\theta[v\circ\pi]=Q_{\theta_B}[v]\circ\pi.
\]
Nevertheless, every $u\in\PSH(X,\theta)$ satisfies
\[
 \mathrm{MA}^{\mathrm{ext}}_\theta(u)=0.
\]
\end{proposition}
\begin{proof}
The form $\theta_B$ is closed, and $\pi$ is a proper holomorphic submersion with connected fibers. Thus \cref{prop:psh-descent-submersion} gives the asserted bijection. Surjectivity of $\pi$ also shows that an ordinary singularity inequality holds upstairs exactly when it holds downstairs, with the same constant.

By \cref{lem:submersion-multiplier-pullback}, for every $s>0$,
\[
 \mathcal I(sv\circ\pi)=\mathcal I(sv)\mathcal O_X.
\]
Extension by a product submersion preserves and reflects inclusion of ideals. To see reflection, compose local generators with a local holomorphic section: a pulled-back germ in the extended ideal restricts to a germ in the original ideal. Thus the all-weight ideal orders and equivalences correspond exactly.

The competitors defining both envelopes correspond under this bijection, including the bound by zero and the required ordinary or ideal condition. Their pointwise suprema commute with pullback. Upper semicontinuous regularization also commutes with pullback in a product chart, since a base neighborhood is obtained from every sufficiently small product neighborhood and local sections give the reverse comparison. This proves the two envelope identities.

Finally, $u_j=\max\{u,-j\}$ descends to $v_j=\max\{v,-j\}$. The semipositivity of $\theta_B$ makes these admissible bounded potentials. On a product chart add a local smooth potential for $\theta_B$ and approximate the resulting bounded plurisubharmonic function on the base by smooth plurisubharmonic functions decreasing on smaller charts. Their pulled-back complex Hessians have at most two complex directions, so their cubes vanish. Bedford--Taylor monotone continuity gives $(\theta+\ddc u_j)^3=0$. The increasing truncation definition of the extended non-pluripolar measure now proves the last assertion.
\end{proof}

Thus equality of the two envelopes in these backgrounds is exactly the corresponding equality on the base, while every three-dimensional Monge--Amp\`ere measure is zero. A positive-volume criterion from the big closed theory cannot be applied to these non-big classes. The strictly Hermitian background $\omega$ belongs to a different theory: its smooth positive-density equations are solvable, even though its infimum of masses is zero.

\subsection{Selected sections and their growth}

\begin{proposition}\label{prop:elliptic-selected-section-descent}
Let $a,b$ be positive integers, let $L_B=\mathcal O_B(a,b)$ and $L=\pi^*L_B$. For every quasi-plurisubharmonic $v$ on $B$ and positive integer $m$, pullback gives an isomorphism
\[
 H^0\bigl(B,L_B^m\otimes\mathcal I(mv)\bigr)
   \longrightarrow
 H^0\bigl(X,L^m\otimes\mathcal I(mv\circ\pi)\bigr).
\]
For $v=0$ these spaces have dimension $(am+1)(bm+1)$.

Equip $L_B$ with a smooth Hermitian metric, choose a continuous real function $\varphi$ on $B$, and put $\nu=\pi_*\mathrm{d}V_X$ for a smooth positive volume form on $X$. Use the norms
\[
 \|s\|_B^2=\int_B|s|^2\mathrm{e}^{-m\varphi}\,\mathrm{d}\nu,
 \qquad
 \|\pi^*s\|_X^2=\int_X|\pi^*s|^2
                \mathrm{e}^{-m\varphi\circ\pi}\,\mathrm{d}V_X.
\]
The isomorphism is isometric. If $K_{m,B}$ and $K_{m,X}$ are the corresponding pointwise squared-norm sums over orthonormal bases, including the displayed continuous weight, then
\[
 K_{m,X}=K_{m,B}\circ\pi,
 \qquad \pi_*(K_{m,X}\,\mathrm{d}V_X)=K_{m,B}\,\mathrm{d}\nu.
\]
\end{proposition}
\begin{proof}
The selected-space identification is \cref{prop:partial-kernel-submersion}, applied to $L_B$ and $v$ on the base. For $v=0$, the usual bihomogeneous polynomial basis on $\mathbf P^1\times\mathbf P^1$ has $(am+1)(bm+1)$ elements.

Integration along the fibers of the proper submersion shows that $\nu$ is a smooth positive volume form. Its defining integration identity proves equality of the two norms. Pull back an orthonormal basis and compare its pointwise norms to obtain the first kernel identity; applying the same integration identity proves the second. The assertions include a zero-dimensional selected space, for which both kernels vanish.
\end{proof}

Even without a singular selecting weight, the growth is quadratic:
\[
 \lim_{m\to\infty}m^{-2}h^0(X,L^m)=ab,
 \qquad \lim_{m\to\infty}m^{-3}h^0(X,L^m)=0.
\]
The exact identities retain the base singularity information and the base kernel problem. A Hermitian metric on the threefold does not change this algebraic growth into cubic growth. This model links the closed-class cone, the descent of singularities and the lower-dimensional kernel theory, while leaving the nonlinear Hermitian equation distinct.


\subsection{The boundary volume and section growth}
\label{sec:elliptic-boundary-volume}

For the fixed balanced volume on this threefold, the zero potential solves the equation for every positive linear combination of $\alpha,\beta,D$. The variational volume is therefore an explicit cubic polynomial. Its order at the nef boundary records which of the two horizontal coefficients vanish.

\begin{proposition}\label{prop:elliptic-boundary-volume}
On the elliptic threefold, put $\theta=a\alpha+b\beta$ with $a,b\geq0$, and let $\mu_0=\omega^3/(24\pi)$. Then for every $\epsilon>0$,
\begin{equation}\label{eq:elliptic-boundary-volume}
 \mathcal V_{\mu_0}(\theta+\epsilon\omega)
       =24\pi\epsilon(a+\epsilon)(b+\epsilon).
\end{equation}
In particular $\mathcal V_\mu(\theta)=0$ for every smooth positive probability measure $\mu$. If $q$ is the number of positive entries among $a,b$, then the orders of \cref{prop:nef-degeneration-orders} satisfy
\[
 \underline a(\theta)=\overline a(\theta)=3-q.
\]
These orders hold for every smooth positive probability measure and every Hermitian perturbation direction.
\end{proposition}
\begin{proof}
The forms $\alpha,\beta,D$ are semipositive of rank one, their squares vanish, and $\omega=\alpha+\beta+D$ is strictly positive. Thus
\[
 \Omega_\epsilon=\theta+\epsilon\omega
       =(a+\epsilon)\alpha+(b+\epsilon)\beta+\epsilon D>0.
\]
The wedge expansion gives
\[
 \Omega_\epsilon^3
   =6\epsilon(a+\epsilon)(b+\epsilon)\alpha\wedge\beta\wedge D
   =24\pi\epsilon(a+\epsilon)(b+\epsilon)\mu_0.
\]
The zero potential therefore solves its smooth normalized equation. By \cref{thm:gaud-volume-dual-formula}, the equation's scalar is precisely $\mathcal V_{\mu_0}(\Omega_\epsilon)$, proving \eqref{eq:elliptic-boundary-volume}. Nef continuity in \cref{thm:gaud-volume-nef-limit} and smooth-density comparison give $\mathcal V_\mu(\theta)=0$.

If both $a,b$ are positive, the leading term is $24\pi ab\epsilon$. If exactly one is positive, it is $24\pi(a+b)\epsilon^2$. If both are zero, the formula is $24\pi\epsilon^3$. The logarithmic quotients have limits $1,2,3$, respectively. Their independence of the density and perturbing metric is \cref{prop:nef-degeneration-orders}.
\end{proof}

\Needspace{9\baselineskip}
For nonnegative integers $a,b$, the bundle $L=\pi^*\mathcal O_B(a,b)$ has
\[
 h^0(X,L^m)=(am+1)(bm+1).
\]
Indeed the compact-fiber descent and bihomogeneous polynomial argument of \cref{prop:elliptic-selected-section-descent} also apply when one or both integers vanish. Its growth degree is exactly $q$, while the boundary order of its smooth curvature representative is $3-q$. Thus the boundary order records the loss of algebraic growth dimension in this explicit family, even though the unperturbed three-dimensional Monge--Amp\`ere measures vanish.

Every positive perturbation $\Omega_\epsilon$ still has zero lower mass, by \cref{thm:vanishing-lower-mass}, and finite upper mass, since $X$ has BMP. Formula \eqref{eq:elliptic-boundary-volume} concerns one fixed smooth density, and does not assert conservation of the raw mass over all potentials. Whether comparable identifications hold beyond such explicit families is part of \cref{qu:nef-degeneration-exponent}.

\subsection{Limits of the minimizing Gauduchon forms}
\label{sec:elliptic-gauduchon-limits}

The minimizing tests for the preceding boundary formula are Gauduchon powers of the equation metrics themselves. Their normalized $(2,2)$-forms converge smoothly, although their metric roots may diverge. The root volume $\lambda_{\mu_0}$ and denominator $D_{\mu_0}$ are those of \cref{def:gaud-variational-volume}; $\asymp$ denotes two-sided bounds by fixed positive multiples as $\epsilon\downarrow0$.

\begin{proposition}\label{prop:elliptic-gauduchon-limits}
Let $a,b\geq0$, $\theta=a\alpha+b\beta$ and $\Omega_\epsilon=\theta+\epsilon\omega$. For every $\epsilon>0$, the metric $\Omega_\epsilon$ is Gauduchon. It minimizes the quotient defining $\lambda_{\mu_0}(\Omega_\epsilon)$. Put
\[
 I_\epsilon=8\pi\bigl(ab+2(a+b)\epsilon+3\epsilon^2\bigr),
 \qquad T_\epsilon=\frac{\Omega_\epsilon^2}{I_\epsilon},
 \qquad \gamma_\epsilon=\frac{\Omega_\epsilon}{\sqrt{I_\epsilon}}.
\]
Then $\int_X\omega\wedge T_\epsilon=1$, $\ddc T_\epsilon=0$ and $T_\epsilon=\gamma_\epsilon^2$. As $\epsilon\downarrow0$, the forms $T_\epsilon$ converge smoothly to
\[
 T_0=\begin{cases}
       (4\pi)^{-1}\alpha\wedge\beta,&a,b>0,\\
       (8\pi)^{-1}\alpha\wedge(\beta+D),&a>0,\ b=0,\\
       (8\pi)^{-1}\beta\wedge(\alpha+D),&a=0,\ b>0,\\
       (24\pi)^{-1}\omega^2,&a=b=0.
      \end{cases}
\]
In every case $\theta\wedge T_0=0$. If $q$ is the number of positive coefficients among $a,b$, the denominator of the normalized minimizing test satisfies
\[
 D_{\mu_0}(\gamma_\epsilon)
  =\frac{\bigl(24\pi\epsilon(a+\epsilon)(b+\epsilon)\bigr)^{2/3}}
         {I_\epsilon}
  \asymp\epsilon^{q/3},
\]
where the two comparison constants may depend on the fixed pair $(a,b)$.
\end{proposition}
\begin{proof}
Write $x=a+\epsilon$, $y=b+\epsilon$ and $z=\epsilon$. Since $\alpha^2=\beta^2=D^2=0$,
\[
 \Omega_\epsilon^2
   =2xy\alpha\wedge\beta+2xz\alpha\wedge D+2yz\beta\wedge D.
\]
The forms $\alpha,\beta$ are closed and $\ddc D=2\alpha\wedge\beta$. Applying $\ddc$ to each term makes its last two terms zero because they contain $\alpha^2$ or $\beta^2$. Thus $\ddc(\Omega_\epsilon^2)=0$. The zero potential solves the normalized equation for $\Omega_\epsilon$ and $\mu_0$, so the characterization of minimizing metrics in \cref{thm:gaud-volume-dual-formula} gives the minimizing assertion.

Fiber integration gives $\int_X\alpha\wedge\beta\wedge D=4\pi$. Hence
\[
 \int_X\omega\wedge\Omega_\epsilon^2
   =8\pi(xy+xz+yz)=I_\epsilon.
\]
This proves the trace normalization. Dividing the three coefficients in the preceding wedge expansion by $I_\epsilon$ gives each displayed limit, with convergence in every smooth norm since the forms are fixed. The identities $\alpha^2=\beta^2=0$ then give $\theta\wedge T_0=0$ in all four cases.

Set $c_\epsilon=24\pi\epsilon(a+\epsilon)(b+\epsilon)$. The equation gives
\[
 \gamma_\epsilon^3
    =\frac{c_\epsilon}{I_\epsilon^{3/2}}\mu_0.
\]
Substitution in the definition of $D_{\mu_0}$ and $\int_X\mu_0=1$ yield $D_{\mu_0}(\gamma_\epsilon)=c_\epsilon^{2/3}/I_\epsilon$. For fixed $(a,b)$ one has
$c_\epsilon\asymp\epsilon^{3-q}$ and $I_\epsilon\asymp\epsilon^{2-q}$. The quotient therefore has order $q/3$, as asserted.
\end{proof}

The limiting positive $(2,2)$-forms have coefficient-matrix ranks $1,2,3$ when $q=2,1,0$, respectively. This follows in a coframe adapted to the three rank-one forms $\alpha,\beta,D$ from the three complementary products in the proof. If $a,b>0$, the metric roots themselves converge smoothly to
\[
 \frac{a\alpha+b\beta}{\sqrt{8\pi ab}}.
\]
If $a>0$ and $b=0$, however, the coefficient of $\alpha$ in $\gamma_\epsilon$ grows as a positive constant times $\epsilon^{-1/2}$, while its other two coefficients tend to zero. Thus the roots are unbounded relative to the fixed metric $\omega$, although their squares converge smoothly. The same statement holds with $\alpha$ and $\beta$ interchanged. A rank-two positive $(2,2)$-form cannot be the square of a finite semipositive $(1,1)$-form in dimension three: the latter's eigenvalues give complementary pairwise products, whose possible nonzero ranks are three or one. At the vertex $a=b=0$, the normalized root instead stays equal to $\omega/\sqrt{24\pi}$.

For a nonzero boundary class, the smooth limiting test annihilates $\theta$ and its determinant denominator vanishes. No concentration of its coefficient measures occurs in this family. This degeneration of fixed-density minimizing tests is different from the sequence of Gauduchon forms used to prove zero lower mass in \cref{lem:lower-mass-unbounded-root}. At the vertex, the denominator remains positive; the order three comes entirely from scaling the background.


\section{An incomplete metric in a big class}
\label{sec:incomplete-divisorial-metric}

We now consider exponential equations $\langle(\theta+\ddc\rho)^n\rangle=\mathrm{e}^{\rho}|\sigma|^{-2}\,\mathrm{d}V$ whose measure has a nonintegrable pole along a divisor, and ask whether the metric $\theta+\ddc\rho$ is complete on the locus where it is smooth and positive. A divisorial part of a solution with minimal singularities can cancel the pole and leave the deleted divisor at finite distance. On the blowup of the projective plane at a point, the resulting metric away from the exceptional divisor is the pullback of the Fubini--Study metric. The class is big but not nef, and the equation has a unique solution of finite energy.

Finite energy refers to the classical closed theory of \cite{SGPT}. For a smooth closed real $(1,1)$-form $\theta$ with big class on a compact K\"ahler surface $Y$, and for $u\in\PSH(Y,\theta)$ with minimal singularities, the Monge--Amp\`ere energy relative to the envelope $V_\theta$ is
\[
 E_\theta(u)=\frac13\sum_{j=0}^2
 \int_Y(u-V_\theta)
 \left\langle(\theta+\ddc u)^j\wedge
                (\theta+\ddc V_\theta)^{2-j}\right\rangle.
\]
For general $u\in\PSH(Y,\theta)$, $E_\theta(u)$ is the decreasing limit of $E_\theta(\max(u,V_\theta-k))$ as $k\to\infty$, and $\mathcal E^1(Y,\theta)$ is the set of potentials with $E_\theta(u)>-\infty$; they have full mass.

\Needspace{18\baselineskip}
\begin{theorem}\label{thm:incomplete-divisorial-metric}\index{Metric completeness!divisorial counterexample}
There exist a smooth projective surface $Y$, a smooth irreducible divisor $E$, a smooth closed real $(1,1)$-form $\theta$ with big class $\alpha$, a Hermitian metric on $\mathcal O_Y(E)$, a defining section $\sigma$ of $\mathcal O_Y(E)$, and a smooth positive volume form $\mathrm{d}V_Y$ with the following properties. The equation
\begin{equation}\label{eq:incomplete-divisorial-equation}
 \left\langle(\theta+\ddc\rho)^2\right\rangle
 =\mathrm{e}^{\rho}|\sigma|^{-2}\,\mathrm{d}V_Y
\end{equation}
has a unique solution $\rho\in\mathcal E^1(Y,\theta)$, and this solution has minimal singularities. Moreover,
\[
 \operatorname{Amp}(\alpha)=Y\setminus E
\]
is precisely the open set on which $\theta+\ddc\rho$ is smooth and strictly positive. The restriction of this metric to every nonempty connected open subset of $Y\setminus E$ is incomplete. In particular, no dense Zariski open subset of $Y\setminus E$ carries a complete induced metric.
\end{theorem}

Let $p=[1:0:0]\in\mathbb P^2$, and let $\omega_2$ be the Fubini--Study form on $\mathbb P^2$, normalized by $\int_{\mathbb P^2}\omega_2^2=1$.

\begin{proof}[Proof of \cref{thm:incomplete-divisorial-metric}]
\emph{The surface and the divisor.}
Use the incidence model
\[
 Y=\{([Z_0:Z_1:Z_2],[W_1:W_2])\in\mathbb P^2\times\mathbb P^1:
           Z_1W_2=Z_2W_1\}.
\]
The projections $\pi\colon Y\to\mathbb P^2$ and $q\colon Y\to\mathbb P^1$ identify $Y$ with the blowup of $\mathbb P^2$ at $p$, with exceptional divisor $E=\pi^{-1}(p)$. Let $\omega_1$ be the Fubini--Study form on $\mathbb P^1$ with the same normalization as $\omega_2$, and write
\[
 \beta=\pi^*\omega_2,\qquad \gamma=q^*\omega_1,
 \qquad \Omega=\beta+\gamma.
\]
The form $\Omega$ is K\"ahler, being the restriction of a product K\"ahler metric.

The linear map $(Z_0,Z_1,Z_2)\mapsto(Z_1,Z_2)$ defines a section $\sigma$ of
\[
 \pi^*\mathcal O_{\mathbb P^2}(1)\otimes q^*\mathcal O_{\mathbb P^1}(-1),
\]
because on $Y$ the vector $(Z_1,Z_2)$ lies in the tautological line $[W_1:W_2]$. The section vanishes simply along $E$, which identifies this line bundle with $\mathcal O_Y(E)$. For the product of the Fubini--Study metric and the metric induced on the tautological bundle, its squared norm is
\[
 s=|\sigma|^2=
 \frac{|Z_1|^2+|Z_2|^2}{|Z_0|^2+|Z_1|^2+|Z_2|^2}.
\]
In a blowup chart $\pi(t,w)=(t,tw)$, $q(t,w)=[1:w]$, one has
\begin{equation}\label{eq:incomplete-divisor-norm}
 s=\frac{|t|^2(1+|w|^2)}{1+|t|^2(1+|w|^2)},
 \qquad \ddc\log s=[E]+\gamma-\beta.
\end{equation}
The current identity follows either from this formula and the second blowup chart or from the Poincar\'e--Lelong formula.

\emph{The big class and its equation.}
Set
\begin{equation}\label{eq:incomplete-current}
 \theta=3\beta-2\gamma,\qquad \rho=2\log s,
 \qquad T=\theta+\ddc\rho=\beta+2[E].
\end{equation}
Thus $\rho\in\PSH(Y,\theta)$. The class $\alpha=\{\theta\}$ is big, because the potential $\psi=(5/2)\log s$ satisfies
\begin{equation}\label{eq:incomplete-strict-current}
 \theta+\ddc\psi
 =\tfrac12(\beta+\gamma)+\tfrac52[E]\geq\tfrac12\Omega.
\end{equation}
It is not nef, since $\int_E\theta=-2$: the restriction of $\beta$ to $E$ vanishes and $q|_E$ is an isomorphism onto $\mathbb P^1$. A nef class has nonnegative degree on every curve, by its defining smooth approximations.

On $Y\setminus E$ define $\mathrm{d}V_Y=\beta^2/s$. To check that it extends across $E$, put
\[
 \mathrm{d}\lambda_{t,w}
 =\tfrac{\mathrm{i}}2\mathrm{d}t\wedge\mathrm{d}\bar t
       \wedge\tfrac{\mathrm{i}}2\mathrm{d}w\wedge\mathrm{d}\bar w,
\]
and let $\mathrm{d}\lambda_{z_1,z_2}$ be the analogous Euclidean volume in affine coordinates of $\mathbb P^2$. The normalized Fubini--Study volume is
\[
 \omega_2^2=\frac{2}{\pi^2}
 \frac{\mathrm{d}\lambda_{z_1,z_2}}{(1+|z_1|^2+|z_2|^2)^3}.
\]
Since the complex Jacobian determinant of $(t,w)\mapsto(t,tw)$ is $t$, we obtain
\begin{align*}
 \beta^2&=\frac{2}{\pi^2}
 \frac{|t|^2\,\mathrm{d}\lambda_{t,w}}
      {[1+|t|^2(1+|w|^2)]^3},\\
 \mathrm{d}V_Y&=\frac{2}{\pi^2}
 \frac{\mathrm{d}\lambda_{t,w}}
      {(1+|w|^2)[1+|t|^2(1+|w|^2)]^2}.
\end{align*}
The second expression is smooth and strictly positive at $t=0$. Interchanging the two homogeneous coordinates gives the extension on the other chart. The two extensions agree off $E$, hence everywhere on their overlap.

The non-pluripolar product of $T$ agrees with $\beta^2$ off $E$ and gives no mass to $E$. The smooth measure $\beta^2$ also gives no mass to $E$, so
\[
 \langle T^2\rangle=\beta^2,
 \qquad
 \mathrm{e}^{\rho}|\sigma|^{-2}\,\mathrm{d}V_Y
   =s^2s^{-1}\frac{\beta^2}{s}=\beta^2.
\]
The right side of \eqref{eq:incomplete-divisorial-equation} is defined on $Y\setminus E$ and extended by zero on $E$, so no product of zero and infinity occurs. This proves the equation and gives $\int_Y\langle T^2\rangle=\int_{\mathbb P^2}\omega_2^2=1$.

\emph{Minimal singularities and energy.}
Let $u\in\PSH(Y,\theta)$. Since
\[
 \beta+\ddc(u-\rho)=\theta+\ddc u\geq0
 \quad\text{on }Y\setminus E,
\]
the function $u-\rho$ on $Y\setminus E$ descends to an $\omega_2$-plurisubharmonic function on $\mathbb P^2\setminus\{p\}$. By \cref{lem:isolated-puncture-extension}, it is bounded above, so $u\leq\rho+C_u$ there. This inequality holds almost everywhere on $Y$, and \cref{lem:qpsh-almost-everywhere-order} makes it hold everywhere. In particular $u=-\infty$ on $E$, and $\rho$ has minimal singularities.

Let $V_\theta$ be the normalized envelope of \cref{sec:obstacle-envelopes}. Since $s\leq1$, we have $\rho\leq0$, and minimality gives
\[
 \rho\leq V_\theta\leq\rho+C.
\]
Thus $\rho-V_\theta$ is bounded off $E$, and the energy is
\[
 E_\theta(\rho)=\frac13\sum_{j=0}^2
 \int_Y(\rho-V_\theta)
 \left\langle(\theta+\ddc\rho)^j\wedge
                (\theta+\ddc V_\theta)^{2-j}\right\rangle.
\]
The mixed measures have finite mass in the closed K\"ahler theory and do not charge $E$. Hence $E_\theta(\rho)>-\infty$, so $\rho\in\mathcal E^1(Y,\theta)$. A potential with minimal singularities has full mass; in this example that mass is $1$.

\emph{Uniqueness.}
Suppose $u\in\mathcal E^1(Y,\theta)$ is another solution, and let $\mathrm{d}\mu=s^{-1}\mathrm{d}V_Y$ on $Y\setminus E$, extended by zero on $E$. This measure is smooth and positive off $E$, but has infinite total mass. The classical comparison principle for full-mass potentials in the closed theory gives
\[
 \int_{\{u<\rho\}}\mathrm{e}^{\rho}\,\mathrm{d}\mu
 \leq\int_{\{u<\rho\}}\mathrm{e}^{u}\,\mathrm{d}\mu.
\]
Both integrals are finite, since the two Monge--Amp\`ere measures have mass $1$. The integrand on the left is strictly larger, which forces $\mu(\{u<\rho\})=0$. Interchanging $u$ and $\rho$ gives $\mu(\{\rho<u\})=0$ as well. Therefore $u=\rho$ almost everywhere on $Y\setminus E$ with respect to a smooth volume form, hence almost everywhere on $Y$, since $E$ has zero volume. By \cref{lem:qpsh-almost-everywhere-order}, $u=\rho$ everywhere.

\emph{The ample locus and incompleteness.}
The K\"ahler current in \eqref{eq:incomplete-strict-current} is smooth off $E$. On the other hand, every $u\in\PSH(Y,\theta)$ satisfies $u\leq\rho+C_u$, so its Lelong numbers at points of $E$ are at least those of $\rho=2\log s$, which are positive. Thus every K\"ahler current in $\alpha$ is singular at every point of $E$, and \cref{def:bc-non-kahler-locus} gives
\[
 \operatorname{Amp}(\alpha)=Y\setminus E_{\mathrm{nK}}(\alpha)
 =Y\setminus E.
\]
The current $T$ equals $\beta$ there, and it is strictly positive because $\pi$ is a biholomorphism away from $E$. Its nonzero divisorial part $2[E]$ prevents smoothness at every point of $E$.

Let $U\subset Y\setminus E$ be a nonempty connected open set, and set $W=\pi(U)$. The map $\pi\colon(U,T|_U)\to(W,\omega_2|_W)$ is an isometry. Choose $x\in W$ and a smooth path $c\colon[0,1]\to\mathbb P^2$ from $x$ to $p$. Its first exit time
\[
 t_0=\min\{t\in[0,1]:c(t)\notin W\}
\]
is positive, and $c([0,t_0))\subset W$. The path has finite Fubini--Study length. For $t_k\nearrow t_0$, the lengths of its tails show that $c(t_k)$ is a Cauchy sequence for the intrinsic metric of $W$. It has no limit in $W$: intrinsic convergence implies ambient convergence, and the ambient limit is $c(t_0)\notin W$. Thus $U$ is incomplete. A dense Zariski open subset of $Y\setminus E$ is the complement of a proper analytic subset in a connected manifold, hence connected, so the last assertion is a special case.
\end{proof}

The cancellation $\mathrm{e}^{\rho}|\sigma|^{-2}=s$ removes the pole, and the divisor $E$, which the metric collapses to the point $p$, stays at finite distance. Thus an exponential equation in a big class can have an incomplete metric even when its finite-energy solution is unique and has minimal singularities. In this example the class is not nef. The next example shows that nefness alone does not restore completeness; there the pole produces a complete cusp, but a different curve lies at finite distance.

\subsection{A nef class with an incomplete null end}
\label{sec:nef-incomplete-cusp}

Nefness does not prevent the failure of completeness. The next construction keeps a complete cusp at the prescribed pole, while a disjoint null curve lies at finite distance.

\begin{theorem}\label{thm:nef-incomplete-cusp}\index{Metric completeness!nef big counterexample}
There are a smooth projective surface $Y$, disjoint smooth divisors $E,D$, a smooth closed semipositive form $\theta$ with nef and big class, a smooth Hermitian norm on $\mathcal O_Y(D)$, its defining section $\sigma$, and a smooth strictly positive volume form $\mathrm{d}V_Y$, such that
\begin{equation}\label{eq:nef-incomplete-cusp-equation}
 \left\langle(\theta+\ddc\rho)^2\right\rangle
 =\mathrm{e}^{\rho}|\sigma|^{-2}\,\mathrm{d}V_Y
\end{equation}
has a unique full-mass solution $\rho\in\mathcal E^1(Y,\theta)$. The solution has zero Lelong numbers and is smooth on $Y\setminus(E\cup D)$, where its current $\Omega=\theta+\ddc\rho$ is strictly positive. The curve $E$ is null for $\{\theta\}$. Every dense Zariski open set on which $\Omega$ is smooth and strictly positive is incomplete, although every path approaching $D$ in $Y\setminus(E\cup D)$ has infinite length.
\end{theorem}

The measure $|\sigma|^{-2}\mathrm{d}V_Y$ has a nonintegrable divisorial pole with smooth strictly positive coefficient. The solution is unbounded at $D$ and does not have ordinary minimal singularities: semipositivity makes the normalized minimal envelope equal to zero. The theorem concerns nefness and bigness; the additional identities required by geometric smoothing remain a separate question.

\begin{lemma}\label{lem:nef-cusp-radial-profile}
There is $F\in C^\infty(\R)$ with $0<p:=F'<1$ and $p'>0$ such that, for $c=1/100$, some $B\in\R$ and $K=(9c^2/8)\mathrm{e}^{-B}>0$,
\begin{align}
 F(t)&=B-3\log(1-c\mathrm{e}^{t/2}),&
 pp'&=K\mathrm{e}^{F+t} &&(t\leq-2),\label{eq:nef-cusp-left-profile}\\
 F(t)&=t+\log H(p(t)),&
 pp'&=\mathrm{e}^{F-t} &&(t\geq2),\label{eq:nef-cusp-right-profile}
\end{align}
where $H(p)=(1-p)^2(1+2p)/6$. At the two ends,
\begin{align*}
 F(t)&=B+3c\mathrm{e}^{t/2}+O(\mathrm{e}^t),&
 p'(t)&\sim(3c/4)\mathrm{e}^{t/2} &&(t\to-\infty),\\
 1-p(t)&\sim2/t,&p'(t)&\sim2/t^2 &&(t\to+\infty),
\end{align*}
and $F(t)-t=-2\log t+\log2+o(1)$ as $t\to+\infty$.
\end{lemma}

\begin{proof}
Put $L(t)=-3\log(1-c\mathrm{e}^{t/2})$ near $(-\infty,-2]$. With $q=c\mathrm{e}^{t/2}$,
\[
 p_L=L'=\frac{3q}{2(1-q)},\qquad
 p_L'=\frac{3q}{4(1-q)^2},\qquad
 p_Lp_L'=\frac{9c^2}{8}\mathrm{e}^{t}(1-c\mathrm{e}^{t/2})^{-3}.
\]
This proves the left identity for every additive constant $B$.

Since $H'(p)=(p-1)p$, the increasing function
\[
 \mathcal T(p)=2+\int_{3/4}^{p}\frac{q}{H(q)}\,\mathrm{d}q
\]
has an inverse $p_R(t)$ for $t\geq2$, extending across $2$, with $p_R(2)=3/4$ and $p_R'=H(p_R)/p_R$. Indeed, its integrand is asymptotic to $2/(1-p)^2$ at $1$. The function $F_R=t+\log H(p_R)$ has derivative $p_R$ and satisfies the right identity. For $q=1-p_R$,
\[
 (1/q)'=\frac{3-2q}{6(1-q)}\longrightarrow\frac12.
\]
Hence $q\sim2/t$, $H(p_R)\sim2/t^2$, and the right-end asymptotics follow.

It remains to interpolate the derivatives. Set $d=p_R(2)-p_L(-2)>0$. Choose disjoint short endpoint collars in $[-2,2]$ and cutoffs $\alpha_L,\alpha_R$ equal to one on smaller collars. For small $\varepsilon>0$,
\[
 Q_0=\alpha_Lp_L'+\alpha_Rp_R'
       +\varepsilon(1-\alpha_L-\alpha_R)
\]
is positive and agrees with the endpoint derivative germs. The collars and $\varepsilon$ can be chosen so that $\int_{-2}^2Q_0<d$. Add a nonnegative smooth interior bump with integral $d-\int_{-2}^2Q_0$. The resulting $Q>0$ gives
\[
 p(t)=p_L(-2)+\int_{-2}^tQ(x)\,\mathrm{d}x.
\]
This matches both endpoint germs, is increasing, and stays in $(0,1)$. Integrating $F'=p$ with $F(2)=F_R(2)$ gives the required $F$ and fixes the left constant $B$. The left expansion follows directly.
\end{proof}

\begin{proof}[Proof of \cref{thm:nef-incomplete-cusp}]
\emph{The background and the pole.}
Use the incidence surface of \cref{thm:incomplete-divisorial-metric}, with blowup projection $\pi:Y\to\mathbb P^2$, second projection $f:Y\to\mathbb P^1$, and exceptional curve $E$ over $[1:0:0]$. Put $D=\pi^{-1}(Z_0=0)$ and $A=\pi^*c_1(\mathcal O(1))$. Then $E\cap D=\varnothing$, $A^2=1$, and $A\cdot E=0$. The class $A$ is nef, since it has a smooth semipositive representative. If $\beta,\gamma$ are the pulled-back Fubini--Study forms, then $\{\gamma\}=A-[E]$ and $\beta+\varepsilon\gamma$ is K\"ahler. Thus
\[
 \frac{\beta+\varepsilon\gamma}{1+\varepsilon}
       +\frac{\varepsilon}{1+\varepsilon}[E]
\]
is a K\"ahler current in $A$, proving bigness.

Choose an even nonnegative smooth function $\zeta$, supported in $[-1,1]$, positive on its interior and of integral one. The smooth convex function
\[
 U(t)=\int_\R\max(0,t-x)\zeta(x)\,\mathrm{d}x
\]
has $U=0$ for $t\leq-1$, $U=t$ for $t\geq1$, and $U',U''\geq0$. On $Z_0\ne0$, write $z=(Z_1/Z_0,Z_2/Z_0)$, $s=|z|^2$, and $t=\log s$. Give the frame $Z_0$ squared norm $\mathrm{e}^{-U(t)}$. It is constant near the affine origin. At infinity, use
\[
 w=Z_0/Z_1,\qquad v=Z_2/Z_1,\qquad h=1+|v|^2,
 \qquad s=h/|w|^2.
\]
The frame $Z_1$ then has squared norm $\mathrm{e}^{-U}/|w|^2=1/h$ near $w=0$. These formulas define a smooth Hermitian metric globally. Pull it back to $\mathcal O_Y(D)$ and let $\sigma=\pi^*Z_0$. Its curvature satisfies
\[
 \theta=\ddc U(t),\quad |\sigma|^2=\mathrm{e}^{-U(t)},\qquad
 \theta=\ddc\log h,\quad |\sigma|^2=|w|^2/h\quad\hbox{near }D.
\]
The radial Hessian proves semipositivity. The form vanishes near $E$ and has rank one near $D$. In particular $V_\theta=0$; the nef class has volume $A^2=1$.

\emph{Extension of the potential.}
On $Y\setminus(E\cup D)\simeq\C^2\setminus\{0\}$, set
\[
 \rho=F(t)-U(t),\qquad \Omega=\ddc F(t).
\]
Its radial Hermitian eigenvalues are $p/s$ and $p'/s$, hence positive. Extend $\rho$ to $E$ by $B$ and to $D$ by $-\infty$.

Near the affine origin the smooth functions $F(\log(|z|^2+\varepsilon))$ are psh and decrease to the continuous extension of $F(\log|z|^2)$. For example, the two Hessian eigenvalues are $p/(s+\varepsilon)$ and $(s p'+\varepsilon p)/(s+\varepsilon)^2$, with $p,p'$ evaluated at $\log(s+\varepsilon)$. Pullback proves positivity near $E$.

Near $D$ put $t_\varepsilon=\log h-\log(|w|^2+\varepsilon)$ and $g(t)=F(t)-t$. The local potentials $g(t_\varepsilon)+\log h$ decrease to $\rho+\log h$. Their Hessians are
\[
 p(t_\varepsilon)\ddc\log h
 +(1-p(t_\varepsilon))\ddc\log(|w|^2+\varepsilon)
 +p'(t_\varepsilon)\frac{\mathrm{i}}{2\pi}
       \partial t_\varepsilon\wedge\bar\partial t_\varepsilon\geq0.
\]
Thus $\rho$ is globally $\theta$-psh. Boundedness at $E$ and the log-log growth at $D$ give zero Lelong numbers everywhere.

\emph{The volume form and full mass.}
With Euclidean area $\mathrm{d}A$, the radial determinant is
\begin{equation}\label{eq:nef-cusp-radial-determinant}
 \Omega^2=\frac{2}{\pi^2}\frac{pp'}{s^2}
                  \mathrm{d}A_{z_1}\mathrm{d}A_{z_2}.
\end{equation}
Define $\mathrm{d}V_Y=\mathrm{e}^{-F}\Omega^2$ off the two divisors. In an exceptional chart $z=(u,uv)$, the real Jacobian is $|u|^2$ and $s=|u|^2h$. The left profile identity gives
\[
 \mathrm{d}V_Y=\frac{2K}{\pi^2h}\,\mathrm{d}A_u\mathrm{d}A_v.
\]
At infinity $z=(1/w,v/w)$ has real Jacobian $|w|^{-6}$. The right identity gives
\[
 \mathrm{d}V_Y=\frac{2}{\pi^2h^3}\,\mathrm{d}A_w\mathrm{d}A_v,
 \qquad
 |\sigma|^{-2}\mathrm{d}V_Y
 =\frac{2}{\pi^2|w|^2h^2}\,\mathrm{d}A_w\mathrm{d}A_v.
\]
The analogous second charts cover each divisor. The extensions are smooth and strictly positive, and agree on overlaps by their equality on the punctured overlaps. The last expression proves the nonintegrable pole with smooth positive coefficient.

On the complement, $\mathrm{e}^{\rho}|\sigma|^{-2}\mathrm{d}V_Y=\Omega^2$. Extend this density by zero on $D$. Both measures give zero mass to the two divisors. Polar integration on $\C^2$ has density $\pi^2s\,\mathrm{d}s$, so
\[
 \int_Y\langle\Omega^2\rangle
 =\int_{-\infty}^{\infty}2pp'\,\mathrm{d}t
 =p(+\infty)^2-p(-\infty)^2=1.
\]
This proves the global equation and full mass.

\emph{Bounded truncations and energy.}
Near $D$ write $\alpha=\ddc t=\ddc\log h$ and $\beta_t=(\mathrm{i}/2\pi)\partial t\wedge\bar\partial t$. They have rank one, and
\[
 \theta=\alpha,\quad \Omega=p\alpha+p'\beta_t,\quad
 \Omega^2=2pp'\alpha\wedge\beta_t,\quad
 \Omega\wedge\theta=p'\alpha\wedge\beta_t,\quad \theta^2=0.
\]
Angular and projective integration of $\alpha\wedge\beta_t$ gives $\mathrm{d}t$. Since $|\rho|=O(\log t)$ and $p'=O(t^{-2})$, its integrals against $\Omega^2$ and $\Omega\wedge\theta$ are finite there. Near $E$ the potential is bounded and $\theta=0$; the remaining region is compact with smooth data.

Choose $M$ so that $\phi=\rho-M\leq0$, and set $\phi_k=\max(\phi,-k)$. For large $k$ the truncation changes only the cusp. There $\phi=\log H(p(t))-M$ is strictly decreasing. If $\phi(T_k)=-k$, then $k=2\log T_k+O(1)$. The punctured radial primitive is
\[
 G_k(t)=\max(F(t)-M,t-k).
\]
Its genuine local potential is $G_k(t)+\log|w|^2$, equal to $\log h-k$ near $D$. This local frame term prevents a spurious divisor mass. The derivative jumps upward from $p(T_k)$ to $1$.

Convolving $G_k$ in $t$ with an even nonnegative smooth kernel gives convex primitives decreasing to $G_k$ as the scale decreases. The local bounded potentials converge uniformly on compact sets, so their Bedford--Taylor products converge. Integrating the smooth radial square and mixed measures through the transition gives respectively the shell masses
\[
 1-p(T_k)^2,\qquad 1-p(T_k).
\]
Beyond the transition the primitives are linear and both top-degree measures vanish. On the shell $\phi_k=-k$, so the absolute energy contributions are $k(1-p(T_k)^2)$ and $k(1-p(T_k))$. Both tend to zero, since $1-p(T_k)\sim2/T_k$. The interior contributions converge absolutely by the preceding bounds. At $E$, bounded potentials have no mass on the curve, consistent with the zero radial boundary mass $p(-\infty)^2=0$. Consequently
\[
 E_\theta(\phi_k)=\frac13\sum_{j=0}^2
       \int_Y\phi_k(\theta+\ddc\phi_k)^j\wedge\theta^{2-j}
\]
has a finite limit and is uniformly bounded below. This proves $\rho\in\mathcal E^1(Y,\theta)$.

Uniqueness follows by the full-mass comparison argument in \cref{thm:incomplete-divisorial-metric}, with $\mathrm{d}\mu=|\sigma|^{-2}\mathrm{d}V_Y$ off $D$, extended by zero on $D$. Its total mass is infinite, but each exponential density is a finite measure of mass $1$. Comparison on the two strict sublevel sets forces equality almost everywhere, and \cref{lem:qpsh-almost-everywhere-order} gives equality of the potentials everywhere.

\emph{The two metric ends.}
Fix $v=v_0$ in an exceptional chart and write $h_0=1+|v_0|^2$. The normal restriction is
\[
 \Omega|_{v=v_0}=\frac{\mathrm{i}}{2\pi}
       \frac{p'(\log(|u|^2h_0))}{|u|^2}
       \,\mathrm{d}u\wedge\mathrm{d}\bar u.
\]
The coefficient is asymptotic to $(3c\sqrt{h_0}/4)|u|^{-1}$, so the current has no smooth extension at any point of $E$. The radial path $u=r>0$ has finite remaining length, bounded by $C\int_0^{r_0}r^{-1/2}\,\mathrm{d}r$. At $D$ the normal coefficient is asymptotic to $2/(|w|^2t^2)$ and also diverges. Thus every smooth positive regular locus excludes both curves.

Let $W$ be any dense Zariski open positive regular locus. The proper analytic set $Y\setminus W$ has finitely many curve and point components, by local finiteness of irreducible components and compactness. Distinct curves meet $E$ in finitely many points. Choose a point of $E$ avoiding all these intersections and point components. In a sufficiently small chart the complement is exactly $E$, so a tail of the preceding radial path lies in $W$. Its points approaching $E$ form an intrinsic Cauchy sequence with no limit in $W$: the smooth positive metric induces its manifold topology, while the only ambient limit is the removed point. Hence $W$ is incomplete.

Finally, on the cusp $\Omega=p\ddc t+p'\,\mathrm{d}t\wedge\dc t$, with both terms nonnegative and $p'\geq c_0/t^2$ for large $t$. The length of any path with $t\to\infty$ is bounded below by a positive multiple of $\int|\mathrm{d}t|/t$, which diverges. The missing finite-distance end is therefore the null curve disjoint from the pole.
\end{proof}


\subsection{An incomplete metric with compatible flow data}
\label{sec:flow-compatible-incomplete-metric}

The auxiliary divisorial equation can be incomplete even when its class, divisor and volume come from the prescribed initial data of a geometric smoothing problem.

Let $(X,\omega)$ be compact K\"ahler, with an $\omega$-psh initial potential $\phi_0$ having analytic singularities. The compatible data on a log resolution $\pi:Y\to X$ satisfy
\begin{gather}
 \pi^*\phi_0=\sum_jm_j\log|\sigma_j|^2+\ell_0,\qquad
 \pi^*(\omega+\ddc\phi_0)=\sum_jm_j[D_j]+R_0,
 \label{eq:compatible-initial-resolution}\\
 K_Y=\pi^*K_X+\sum_jb_jD_j,\qquad
 \pi^*\mathrm{d}V_X=\prod_j|\sigma_j|^{2b_j}\mathrm{d}V_Y.
 \label{eq:compatible-volume-resolution}
\end{gather}
Here $D=\sum_jD_j$ is SNC, the defining sections have smooth Hermitian squared norms, $\ell_0$ is smooth, and $R_0$ is smooth closed semipositive. If $\Theta_j$ denotes the curvature of the section norm, put
\[
 \lambda=\max_j\frac{m_j}{1+b_j},\qquad
 \Theta=\sum_j(1+b_j)\Theta_j,\qquad
 \theta=\lambda\Theta+R_0,\qquad
 |\sigma|^2=\prod_j|\sigma_j|^2.
\]
The number $\lambda$ is the inverse exponential-integrability index of $\phi_0$. The metric in the next theorem is the solution of the auxiliary elliptic equation at this background.

\begin{theorem}\label{thm:flow-compatible-incomplete-metric}\index{Metric completeness!compatible initial data}
There are compatible data \eqref{eq:compatible-initial-resolution}--\eqref{eq:compatible-volume-resolution} on smooth projective surfaces, with $\mathrm{d}V_X=\omega^2$, $R_0$ K\"ahler and $D$ a disjoint union of two smooth curves, such that
\begin{equation}\label{eq:compatible-incomplete-equation}
 \langle(\theta+\ddc\rho)^2\rangle
 =\mathrm{e}^{\rho}|\sigma|^{-2}\,\mathrm{d}V_Y
\end{equation}
has a unique full-mass solution in $\mathcal E^1(Y,\theta)$. It has the Lelong numbers of a minimal potential, but every dense Zariski open positive regular locus of its current is incomplete. The unweighted pole is nonintegrable, with smooth strictly positive coefficient across both components of $D$.
\end{theorem}

The auxiliary class in this example is non-nef. Thus the theorem and \cref{thm:nef-incomplete-cusp} answer the two separate completeness assertions negatively. Completeness under simultaneous nefness and these compatibility identities remains open. Neither example identifies its auxiliary metric with the evolving maximal flow.

\Needspace{12\baselineskip}
\begin{lemma}\label{lem:fixed-volume-curve-cusp}
On $\mathbb P^1$, let $z=Z_1/Z_0$, $r=|z|^2$, $\alpha=\ddc\log(1+r)$ and $h=r/(1+r)$, with $\int\alpha=1$. For every $c>2$ there is a rotational $c\alpha$-psh function $u$, smooth off $0$, such that
\begin{equation}\label{eq:fixed-volume-curve-cusp}
 g_1:=c\alpha+\ddc u=\mathrm{e}^{u}h^{-1}c\alpha.
\end{equation}
It has zero Lelong number at $0$, mass $\int g_1=c$, and
\[
 \int_{\mathbb P^1}|u|\alpha+\int_{\mathbb P^1}|u|g_1<\infty.
\]
For $s=-\log r\to\infty$, $t=\log r$ and some $A\in\R$,
\[
 \begin{aligned}
 u&=\log(2/c)-2\log(s+A)+O(\mathrm{e}^{-s}),\\
 u_t&=\frac{2}{s+A}+O(\mathrm{e}^{-s}),\qquad
 u_{tt}=\frac{2}{(s+A)^2}+O(\mathrm{e}^{-s}).
 \end{aligned}
\]
\end{lemma}

\begin{proof}
\emph{Barriers and the radial inverse.}
Set $L=1-\log h=1+\log(1+1/r)$ and $w=-2\log L$. Differentiation gives
\[
 \ddc w=\left(-\frac2L+\frac2{rL^2}\right)\alpha,\qquad
 Q(r):=\frac{c\alpha+\ddc w}{\mathrm{e}^{w}h^{-1}c\alpha}
 =\frac{crL^2-2rL+2}{c(1+r)}.
\]
The function $Q$ is positive, since $L\geq1$ and $c>2$. Its limits are $2/c$ at $0$ and $1-2/c$ at infinity, so $0<q_-\leq Q\leq q_+<\infty$. Choose $a_0\leq\min(0,\log q_-)$ and $b_0\geq\max(0,\log q_+)$. Then $l=w+a_0$ and $H=w+b_0$ satisfy
\[
 c\alpha+\ddc l\geq\mathrm{e}^{l}h^{-1}c\alpha,\qquad
 c\alpha+\ddc H\leq\mathrm{e}^{H}h^{-1}c\alpha,\qquad l\leq w\leq H.
\]

Exhaust $\mathbb P^1\setminus\{0\}$ by disks $|\zeta|<R$ in $\zeta=1/z$, with boundary value $w$. For $q=|\zeta|^2$, the equation becomes
\[
 \Delta u=f(q)\mathrm{e}^{u}-g(q),\qquad
 f(q)=\frac{4c}{1+q},\qquad g(q)=\frac{4c}{(1+q)^2},
\]
where $\Delta$ is the Euclidean real Laplacian.

For $k>0$, the regular radial solution $v=V(q)$ of $(-\Delta+k)v=F(q)$ satisfies
\begin{equation}\label{eq:radial-cusp-volterra}
 V(q)=A_0+\frac14\int_0^q\log(q/t)\bigl(kV(t)-F(t)\bigr)\,\mathrm{d}t.
\end{equation}
The $j$-fold iterate of the integral kernel on $1$ is $q^j/(4^j(j!)^2)$, so successive integration converges uniformly on every finite interval. The coefficient of $A_0$ is $\sum_{j\geq0}(kq/4)^j/(j!)^2>0$. Hence a unique $A_0$ gives any prescribed value at $q=R^2$. Moreover,
\[
 V'(q)=\frac14\int_0^1\bigl(kV(qt)-F(qt)\bigr)\,\mathrm{d}t.
\]
This formula gives smoothness at $q=0$ by induction. The maximum principle for $-\Delta+k$ proves uniqueness and order preservation of the radial inverse.

Choose $k\geq\sup f\mathrm{e}^{H}$ on the fixed disk. Invert $-\Delta+k$ with boundary value $w$ and right side $kv-f\mathrm{e}^{v}+g$, which is nondecreasing on $[l,H]$. The barrier inequalities give $l\leq T(l)$ and $T(H)\leq H$. Iteration from $l$ therefore increases and stays below $H$. The displayed derivative formula gives a uniform first-derivative bound on the fixed interval. Compactness and monotonicity give uniform convergence; passing to \eqref{eq:radial-cusp-volterra} proves the equation, and the averaged derivative formula bootstraps the limit to smoothness.

The same barriers bound the solutions on every fixed reciprocal disk as $R\to\infty$. For an exact nonlinear solution $U(q)$,
\[
 U'(q)=\frac14\int_0^1\bigl(f(qt)\mathrm{e}^{U(qt)}-g(qt)\bigr)\,\mathrm{d}t.
\]
This gives all derivative bounds on compact intervals by induction. A diagonal limit is smooth on the reciprocal plane, including $z=\infty$, and has $u=-2\log(-\log r)+O(1)$ at $0$.

\emph{The asymptotics, mass and energy.}
The radial equation is
\[
 u_{tt}=\frac{c\mathrm{e}^{u}}{1+\mathrm{e}^t}
              -\frac{c\mathrm{e}^t}{(1+\mathrm{e}^t)^2}.
\]
Put $v(t)=u(t)+c\log(1+\mathrm{e}^t)$ and $V(s)=v(-s)$. Then
\[
 V''=k_0(s)\mathrm{e}^{V},\qquad
 k_0(s)=c(1+\mathrm{e}^{-s})^{-(c+1)},\qquad
 V=-2\log s+O(1).
\]
Convexity and this sublinear decay imply $V'<0$ and $V'\to0$ at infinity. The function $J=(V')^2-2k_0\mathrm{e}^{V}$ satisfies $J'=-2k_0'\mathrm{e}^{V}$ and tends to zero. Since $k_0'=O(\mathrm{e}^{-s})$ and $\mathrm{e}^{V}\asymp s^{-2}$,
\[
 J(s)=2\int_s^\infty k_0'(\tau)\mathrm{e}^{V(\tau)}\,\mathrm{d}\tau
       =O(\mathrm{e}^{-s}s^{-2}).
\]
Thus $(\mathrm{e}^{-V/2})'=\sqrt{c/2}(1+O(\mathrm{e}^{-s}))$, and integration gives $\mathrm{e}^{-V/2}=\sqrt{c/2}(s+A)+O(\mathrm{e}^{-s})$. This identity and the differential equation give all three asserted asymptotics.

The local potential $u+c\log(1+r)$ extends subharmonically by \cref{lem:qpsh-removable-analytic}. Its radial slope tends to zero at the puncture, so the extension adds no point mass. The density is asymptotic to
\[
 \frac{\mathrm{i}}{2\pi}\frac{2}{r s^2}
          \,\mathrm{d}z\wedge\mathrm{d}\bar z.
\]
Distributional Stokes gives mass $c$. The energy integral at the cusp is bounded by $C\int^\infty(\log s)/s^2\,\mathrm{d}s$; the integral against the smooth form $\alpha$ is finite as well.
\end{proof}

\Needspace{12\baselineskip}
\begin{proof}[\cref{thm:flow-compatible-incomplete-metric}]
\emph{The initial data and the resolution.}
Take $X=\mathbb P^1_Z\times\mathbb P^1_W$, with normalized Fubini--Study forms $\alpha,\beta$, and put
\[
 \omega=6(\alpha+\beta),\qquad \mathrm{d}V_X=\omega^2=72\alpha\wedge\beta.
\]
Let $H=(Z_1=0)$ and $p=([0:1],[1:0])=(\infty,0)$. Define
\[
 \begin{gathered}
 h_H=\frac{|Z_1|^2}{|Z_0|^2+|Z_1|^2},\qquad
 \eta=\frac{|Z_0|^2}{|Z_0|^2+|Z_1|^2}
              +\frac{|W_1|^2}{|W_0|^2+|W_1|^2},\\
 \phi_0=2\log h_H+2\log\eta.
 \end{gathered}
\]
The zero sets are $H$ and $\{p\}$, which are disjoint. Blow up $p$, write $E$ for the exceptional curve, and set $H_Y=\pi^{-1}H$, $D=H_Y+E$. The projective blowup is smooth, compact, connected and K\"ahler, and is an isomorphism off the initial singular locus.

Use $|\sigma_H|^2=\pi^*h_H$ and $|\sigma_E|^2=\pi^*\eta$. Near $p$, put $x=Z_0/Z_1$, $y=W_1/W_0$. In the chart $x=a$, $y=ab$,
\[
 \pi^*\eta=|a|^2A(a,b),\qquad
 A(a,b)=\frac1{1+|a|^2}+\frac{|b|^2}{1+|ab|^2}>0.
\]
This defines a smooth Hermitian metric on $\mathcal O_Y(E)$ across $E$. The second chart is obtained by interchanging $x,y$, and the metric factors satisfy the transition rule.

The sections $Z_0W_0$, $\sqrt2Z_0W_1$, $Z_1W_1$ of $\mathcal O(1,1)$ have summed squared norm $\eta$ and base ideal $(x,\sqrt2xy,y)=(x,y)$. Dividing by $\sigma_E$ gives a morphism $F:Y\to\mathbb P^2$ with
\begin{equation}\label{eq:compatible-curvature-map}
 F^*\gamma=\pi^*(\alpha+\beta)-\Theta_E,\qquad
 \Theta_H=\pi^*\alpha,
\end{equation}
where $\gamma$ is normalized Fubini--Study. In the first chart $F=[1:\sqrt2ab:b]$, so $F|_E$ is an isomorphism onto a line.

Near $p$,
\[
 \eta=\frac{|x|^2+|y|^2+2|xy|^2}{(1+|x|^2)(1+|y|^2)}.
\]
Thus $\phi_0$ has analytic singularities; in the convention $\frac1{2m}\log\sum|f_j|^2$, the coefficient $2$ is obtained by taking the fourth power of the sum of squares and expanding it as a sum of squared holomorphic monomials. The exact resolved identities are
\begin{align}\label{eq:compatible-kahler-residual}
 \pi^*\phi_0&=2\log|\sigma_H|^2+2\log|\sigma_E|^2,\\
 R_0&=\pi^*\omega-2\Theta_H-2\Theta_E
       =2\pi^*\alpha+4\pi^*\beta+2F^*\gamma.\notag
\end{align}
The first two terms of $R_0$ are positive away from $E$. On $E$, their kernel is $TE$, and $F^*\gamma$ is positive on $TE$. Hence $R_0$ is K\"ahler. The resolved current is $2[H_Y]+2[E]+R_0$, which proves $\omega$-psh positivity away from $p$; the bounded-above extension across $p$ gives it everywhere. We have $m_H=m_E=2$ and $\ell_0=0$.

The blowup Jacobian is $a$, so $b_E=1$, $b_H=0$. Direct exponential integrability requires $q<1/2$ at $H$ and $q<1$ at $p$. Consequently $\lambda=2=\max\{2,1\}$, and
\begin{equation}\label{eq:compatible-auxiliary-background}
 \Theta=\Theta_H+2\Theta_E,\qquad
 \theta=2\Theta+R_0=\pi^*\omega+2\Theta_E.
\end{equation}
Since $\Theta_E|_E=-\omega_{\mathrm{FS}}$, the class has intersection $-2$ with $E$ and is non-nef. It is big: the current associated to $3\log|\sigma_E|^2$ is
\[
 (\pi^*\omega-\Theta_E)+3[E],\qquad
 \pi^*\omega-\Theta_E=5\pi^*(\alpha+\beta)+F^*\gamma>0.
\]

\emph{The exact prescribed volume.}
Define $\mathrm{d}V_Y=\pi^*\omega^2/|\sigma_E|^2$ off $E$. In the chart above,
\[
 \mathrm{d}V_Y=\frac{72\,\mathrm{d}A_a\mathrm{d}A_b}
 {\pi^2A(a,b)(1+|a|^2)^2(1+|ab|^2)^2}.
\]
The cancellation of the Jacobian $|a|^2$ is exact. The second chart gives the extension on the rest of $E$, and the extensions agree on overlaps. Thus the volume is globally smooth and strictly positive, with $\pi^*\mathrm{d}V_X=|\sigma_E|^2\mathrm{d}V_Y$ as required. Since $H_Y$ and $E$ are disjoint, $|\sigma_H\sigma_E|^{-2}\mathrm{d}V_Y$ has a nonintegrable ordinary pole with smooth positive coefficient at each component.

\emph{The solution and the minimal reference.}
Take $u$ from \cref{lem:fixed-volume-curve-cusp} with $c=6$, pulled to the first factor of $X$, and put $g_X=(6\alpha+\ddc u)+6\beta$. It is smooth and strictly positive near $p$, and
\[
 g_X^2=\mathrm{e}^{u}h_H^{-1}\omega^2.
\]
Set $P=2\log|\sigma_E|^2$ and $\rho=P+\pi^*u$. Then
\[
 \theta+\ddc P=2[E]+\pi^*\omega,\qquad
 \theta+\ddc\rho=2[E]+\pi^*g_X.
\]
Off $D$ the equation follows from the exact identity
\[
 \mathrm{e}^{\rho}|\sigma_H\sigma_E|^{-2}\mathrm{d}V_Y
 =\mathrm{e}^{\pi^*u}|\sigma_H|^{-2}|\sigma_E|^2\mathrm{d}V_Y
 =\pi^*g_X^2.
\]
Both measures extend by zero on $D$, proving the global non-pluripolar equation.

For any $\theta$-psh potential $\psi$, the function $\psi-P$ descends off $E$ to an $\omega$-psh function on $X\setminus\{p\}$. By \cref{lem:isolated-puncture-extension}, it extends and is bounded above on $X$. Thus $\psi\leq P+C_\psi$ off $E$, and \cref{lem:qpsh-almost-everywhere-order} gives the inequality everywhere. Since $P$ itself is $\theta$-psh, it has minimal singularities.

Its current has non-pluripolar mass $\int_X\omega^2=72$, the volume of the class. In contrast, the ordinary cup-product square is $72+4E^2=68$. The cusp has mass $6$ and no atom, so the current at $\rho$ also has mass $\int_Xg_X^2=72$. Thus $\rho$ has full mass. Near $E$, $\pi^*u$ is smooth, giving the Lelong numbers of $P$; on $H_Y$ the log-log cusp has zero Lelong number. The actual solution is not of ordinary minimal type, because $u$ is unbounded below at $H$.

\emph{Relative energy and uniqueness.}
Choose $M$ with $u-M\leq0$ and use the minimal reference $P+M$. Its canonical approximants are
\[
 \rho_k=P+M+\pi^*u_k,\qquad u_k=\max(u-M,-k).
\]
Their relative difference off $E$ is $\pi^*u_k$. On the curve let $b(t)=6\log(1+\mathrm{e}^t)$ and $V(t)=b(t)+u(t)$. For large $k$ there is a unique cusp level $t=T_k$ with $u(T_k)-M=-k$. If $S_k=-T_k$, the asymptotics give $S_k\asymp\mathrm{e}^{k/2}$.

The truncated curve primitive is $b+u_k=\max(V-M,b-k)$. Below $T_k$ its mass is $b'(T_k)$; at $T_k$ the derivative jumps upward by $u_t(T_k)$. Smooth convex approximation and bounded Bedford--Taylor continuity therefore give total cutoff mass $V'(T_k)=O(S_k^{-1})$. On that region $u_k=-k$, so its absolute energy contribution is
\[
 kV'(T_k)=O(k\mathrm{e}^{-k/2})\longrightarrow0.
\]
Outside it the original density and potential are unchanged, and $\int|u|g_1$ is finite. Consequently $\int|u_k|g_{1,k}$ has a finite limit, where $g_{1,k}=6\alpha+\ddc u_k$. Also $\int|u_k|\alpha$ converges by domination.

The mixed measures for the product current and its reference are
\[
 \omega^2=72\alpha\wedge\beta,\quad
 (g_{1,k}+6\beta)\wedge\omega
     =6g_{1,k}\wedge\beta+36\alpha\wedge\beta,\quad
 (g_{1,k}+6\beta)^2=12g_{1,k}\wedge\beta.
\]
Their pullbacks are precisely the mixed non-pluripolar measures upstairs, ignoring $E$. All three canonical relative energy integrals therefore have finite limits. This proves $\rho\in\mathcal E^1(Y,\theta)$. Replacing $P+M$ by the normalized minimal reference changes it by a bounded amount and preserves this finite-energy class.

Uniqueness follows by the same full-mass comparison argument as in \cref{thm:incomplete-divisorial-metric}. The unweighted pole measure has infinite mass, but the two exponential densities being compared have mass $72$. Comparison on the strict sublevel sets forces equality almost everywhere, and then equality of the canonical quasi-psh representatives.

\emph{The missing metric boundary.}
The positive smooth locus of the current is $Y\setminus(H_Y\cup E)$. The cusp coefficient diverges at $H_Y$, and the nonzero current $2[E]$ excludes smoothness at $E$. The pullback metric part is also degenerate along $TE$.

For any dense Zariski open positive regular locus $U$, choose a point of $E$ avoiding every other component of its analytic complement, as in \cref{thm:nef-incomplete-cusp}. In a small chart the complement is exactly $E$. A transverse path $a=\xi>0$, $b=b_0$ lies in $U$ for small $\xi$ and approaches that point. Since $g_X$ is smooth near $p$ and $\pi$ is smooth, the coefficients of $\pi^*g_X$ are bounded in the chart. The path has bounded speed and finite remaining length. Its tails give an intrinsic Cauchy sequence with no limit in $U$, since its only ambient limit is the removed point and a smooth positive Riemannian metric induces the manifold topology. Hence every such $U$ is incomplete.
\end{proof}


\section{Normal length near a divisor}
\label{sec:divisorial-normal-length}

To prove completeness near a divisor, one needs a lower bound for the metric in the normal direction. A lower bound for the determinant alone does not give it, because the tangential eigenvalues may absorb the volume. Dividing the determinant by the determinant of the tangential block gives the normal Schur complement. A lower bound for it by the inverse square of a defining function of the divisor, weakened by at most the square of the logarithm of that function, forces completeness. For solutions of exponential equations, this bound follows from bounds on the density, the potential and the tangential determinant, and an example on a torus shows that the logarithmic exponent $2$ is sharp.

In this section $X$ is again a connected compact complex manifold of dimension $n\geq1$ with a Hermitian form $\omega$. Let $D\subset X$ be a nonempty smooth effective divisor, and put $M=X\setminus D$. Let $s_D$ be the canonical section of $\mathcal O_X(D)$, equipped with a smooth Hermitian norm, and set
\[
 R=|s_D|,\qquad \tau=-\log R^2.
\]
Rescale the norm so that $\tau\geq1$ on $M$. Write $J$ for the complex structure. For a smooth positive real $(1,1)$-form $T$ on $M$, we use the Riemannian metric $g_T(v,v)=\pi T(v,Jv)$ on real tangent vectors $v$. Thus, if $T=(\mathrm{i}/2\pi)\sum H_{j\bar k}\,\mathrm{d}z_j\wedge\mathrm{d}\bar z_k$, then $g_T(v,v)=\sum H_{j\bar k}\,\mathrm{d}z_j(v)\overline{\mathrm{d}z_k(v)}$, and Hermitian norms below use this normalization of coefficients.

\Needspace{7\baselineskip}
Near $D$, the form $\partial\tau$ does not vanish (see \eqref{eq:normal-gradient-comparison} below), and we define the tangential hyperplane field, the tangential block and the determinant
\[
 E_\tau=\ker\partial\tau,\qquad B_\tau=T|_{E_\tau},\qquad
 \det_\omega T=\frac{T^n}{\omega^n}.
\]
Similarly, $\det_\omega B_\tau$ denotes the determinant of $B_\tau$ relative to the restriction of $\omega$ to the complex hyperplane $E_\tau$; when $n=1$, it is $1$.

\begin{proposition}\label{prop:divisorial-normal-length}\index{Metric completeness!normal length}
Suppose that there are constants $c>0$, $\tau_0\geq1$ and $p\leq2$ such that, on $\{\tau>\tau_0\}$,
\begin{equation}\label{eq:normal-determinant-ratio}
 \frac{\det_\omega T}{\det_\omega B_\tau}
 \geq cR^{-2}\tau^{-p}.
\end{equation}
Then $(M,g_T)$ is complete.
\end{proposition}

\begin{proof}
In a coordinate neighborhood meeting $D$, write $s_D=z_1 e$, where $e$ is a local holomorphic frame with $|e|^2=\mathrm{e}^{-\psi}$ and $\psi$ is smooth. Then
\[
 \tau=-\log|z_1|^2+\psi,\qquad
 \partial\tau=-\frac{\mathrm{d}z_1}{z_1}+\partial\psi.
\]
Covering the compact divisor by finitely many such neighborhoods and comparing with the smooth background, we obtain, after increasing $\tau_0$,
\begin{equation}\label{eq:normal-gradient-comparison}
 c_1R^{-2}\leq |\partial\tau|_\omega^2\leq C_1R^{-2}
 \quad\hbox{on }\{\tau>\tau_0\}
\end{equation}
for positive constants $c_1,C_1$. In particular, $E_\tau$ is a hyperplane there.

At a point of this set, choose an $\omega$-unitary frame whose last $n-1$ vectors span $E_\tau$. Let $H$ be the Hermitian matrix that represents $g_T$ as the quadratic form $w^*Hw$ in the complex components $w$, and write
\[
 H=\begin{pmatrix}h&a^*\\a&B\end{pmatrix},\qquad
 S=h-a^*B^{-1}a=\frac{\det H}{\det B},
\]
where $a$ is a column vector and $a^*$ its conjugate transpose. Completing the square gives, for a vector with components $(\xi,\eta)$,
\[
 \begin{pmatrix}\bar\xi&\eta^*\end{pmatrix}
 H\begin{pmatrix}\xi\\\eta\end{pmatrix}
 =S|\xi|^2+(\eta+B^{-1}a\xi)^*B(\eta+B^{-1}a\xi)
 \geq S|\xi|^2.
\]
For $n=1$ the same identity holds without the second term. For the complex components of a real vector $v$ we have
$|\xi|=|\partial\tau(v)|/|\partial\tau|_\omega$, so \eqref{eq:normal-determinant-ratio} and \eqref{eq:normal-gradient-comparison} imply
\begin{equation}\label{eq:normal-length-bound}
 g_T(v,v)\geq c_2\tau^{-p}|\partial\tau(v)|^2
 \geq \frac{c_2}{4}\tau^{-p}|\mathrm{d}\tau(v)|^2,
\end{equation}
where $c_2>0$ is uniform. The last inequality uses $\mathrm{d}\tau(v)=2\operatorname{Re}\partial\tau(v)$.

Define a continuous function $\Phi_p$ on $[1,\infty)$ by
\[
 \Phi_p(t)=
 \begin{cases}
  0,&t\leq\tau_0,\\
  \displaystyle\int_{\tau_0}^t x^{-p/2}\,\mathrm{d}x,&t>\tau_0.
 \end{cases}
\]
Integrating \eqref{eq:normal-length-bound} along piecewise smooth curves shows that $\Phi_p\circ\tau$ is Lipschitz for the length distance of $g_T$; this also holds for curves crossing $\{\tau=\tau_0\}$, since $\Phi_p$ is locally Lipschitz and constant below that level. Since $p\leq2$, $\Phi_p(t)$ tends to infinity with $t$, so $\Phi_p\circ\tau$ is a proper function on $M$.

A Cauchy sequence for the length distance has bounded values of this Lipschitz function and hence lies in a compact subset of $M$. It therefore has a subsequence converging in the manifold topology. On a coordinate ball about the limit, the smooth positive form $T$ is comparable to the Euclidean metric, so this subsequence also converges for the length distance. The Cauchy property then gives convergence of the whole sequence, which proves completeness.
\end{proof}

\Needspace{16\baselineskip}
\begin{corollary}\label{cor:divisorial-exponential-completeness}\index{Exponential equation!normal completeness}
In the preceding setting, let $\theta$ be a smooth real $(1,1)$-form on $X$, and let $u\in C^\infty(M)$ satisfy $T=\theta+\ddc u>0$ on $M$. Suppose that, for some $\lambda>0$ and some positive $f\in C^\infty(M)$,
\[
 T^n=\mathrm{e}^{\lambda u}f\omega^n\quad\hbox{on }M.
\]
Assume that, near $D$, there are constants $c,C>0$ and $q,a,b\in\mathbb R$ such that
\begin{equation}\label{eq:divisorial-three-bounds}
 f\geq cR^{-2}\tau^{-q},\qquad
 u\geq-\frac{a}{\lambda}\log\tau-C,\qquad
 \det_\omega B_\tau\leq C\tau^b.
\end{equation}
If $q+a+b\leq2$, then $(M,g_T)$ is complete.
\end{corollary}

\begin{proof}
The equation and the first two inequalities give
$\det_\omega T\geq c\mathrm{e}^{-\lambda C}R^{-2}\tau^{-(q+a)}$.
Dividing by the third bound gives \eqref{eq:normal-determinant-ratio} with $p=q+a+b$, and \cref{prop:divisorial-normal-length} applies.
\end{proof}

The relevant tangential determinant is that of the intrinsic hyperplane field $E_\tau$, not that of an arbitrary block of coordinate directions.

\subsection{The critical exponent on a compact torus}

The exponent $2$ in \cref{cor:divisorial-exponential-completeness} is sharp. On a product torus, two explicit logarithmic profiles give a complete metric with $q+a+b=2$ and incomplete metrics with $q+a+b=2+\beta$ for every $\beta>0$. Both profiles have full mass, and the second one is bounded.

Let $E=\mathbb C/(\mathbb Z+\mathrm{i}\mathbb Z)$, let $\omega_E=(\mathrm{i}/2\pi)\,\mathrm{d}z\wedge\mathrm{d}\bar z$, and put $X=E^n$, $D=\{0\}\times E^{n-1}$ and $M=X\setminus D$. Write $\pi_j\colon X\to E$ for the coordinate projections and $\omega_0=\sum_{j=1}^n\pi_j^*\omega_E$. Choose radii $0<r_1<r_2<\mathrm{e}^{-1}$ and a smooth radial cutoff $\chi$ on the coordinate disk $\{|z|<r_2\}\subset E$, equal to $1$ on $\{|z|\leq r_1\}$ and to $0$ near its outer boundary, with $0\leq\chi\leq1$. For fixed numbers $A,\eta,\beta>0$, set $L=-\log|z|^2$ and define on $E$
\begin{equation}\label{eq:torus-length-profiles}
 v_{\mathrm c}=-A\chi\log L,\qquad
 v_\beta=\eta\chi L^{-\beta}.
\end{equation}
Both functions are extended by zero outside the disk, and at its center we set $v_{\mathrm c}(0)=-\infty$ and $v_\beta(0)=0$. The bounded profile also occurs in \cref{ex:logarithmic-continuous-density}. We write $\mathrm{MA}_\theta(u)$ for the extended non-pluripolar Monge--Amp\`ere measure $\mathrm{MA}^{\mathrm{ext}}_\theta(u)$ of \cref{thm:arbitrary-extended-ma}.

\Needspace{14\baselineskip}
\begin{lemma}\label{lem:torus-logarithmic-full-mass}
For $C>0$ sufficiently large, put $\theta=C\omega_0$. For either profile $v$ in \eqref{eq:torus-length-profiles}, the function $u=\pi_1^*v$ belongs to $\PSH(X,\theta)$, the form $T=\theta+\ddc u$ is smooth and positive on $X\setminus D$, and
\[
 \int_X\mathrm{MA}_\theta(u)=\int_X\theta^n.
\]
The same $C$ may be used for the two fixed profiles. The bounded profile has minimal ordinary singularities; the profile $v_{\mathrm c}$ does not.
\end{lemma}

\begin{proof}
For a twice differentiable function $G$ of $L=-\log|z|^2$, differentiation away from zero gives
\[
 \frac{\partial^2G(L)}{\partial z\partial\bar z}
 =\frac{G''(L)}{|z|^2}.
\]
On the inner disk this yields
\begin{equation}\label{eq:torus-profile-hessians}
 (v_{\mathrm c})_{z\bar z}=\frac{A}{|z|^2L^2},\qquad
 (v_\beta)_{z\bar z}=\frac{\eta\beta(\beta+1)}{|z|^2L^{\beta+2}}.
\end{equation}
Both densities are integrable at zero. Both profiles are locally integrable, and their distributional Hessians have no additional atom at zero. To see this, integrate twice by parts outside $\{|z|=\varepsilon\}$: the boundary terms against a smooth test function are bounded by constant multiples of
$\varepsilon|G(-\log\varepsilon^2)|$ and $|G'(-\log\varepsilon^2)|$, and both tend to zero for the two profiles. Thus \eqref{eq:torus-profile-hessians} holds in the sense of distributions across the center. On the cutoff annulus all derivatives are bounded. Hence, for a sufficiently large common $C$,
\[
 \alpha_v=C\omega_E+\ddc v
\]
is a positive measure, smooth and strictly positive away from zero, with no atom at zero. Pairing $\ddc v$ with the constant test function gives
$\int_E\alpha_v=C\int_E\omega_E$.

Let $Y=E^{n-1}$, let $\pi_Y\colon X\to Y$ be the projection onto the remaining factors, and let $\theta_Y$ be $C$ times the sum of the $n-1$ flat forms of $Y$. On $X\setminus D$ the smooth wedge product is
\[
 T^n=n\,\pi_1^*\alpha_v\wedge\pi_Y^*\theta_Y^{n-1};
\]
for $n=1$ this means just $\alpha_v$, with $Y$ a point. The right side defines a finite product measure on $X$ that gives $D$ zero mass. By \cref{thm:arbitrary-extended-ma}, the non-pluripolar measure agrees with this product off $D$ and gives the pluripolar set $D$ zero mass, so the two global measures are equal. Integrating the product proves the full-mass identity. Since $\theta>0$, its envelope is $V_\theta=0$ (\cref{sec:obstacle-envelopes}). The bounded profile therefore has minimal ordinary singularities, whereas $v_{\mathrm c}$ is not bounded below by $0$ minus a constant.
\end{proof}

We apply the preceding results with $\omega=\theta$. Choose a smooth norm on $\mathcal O_E(0)$, the line bundle of the point $0\in E$, for which its canonical section has squared norm $|z|^2$ near zero; such a norm is obtained from any smooth norm by multiplying by a smooth positive function prescribed on a smaller disk. Pull it back to $X$ and rescale it by a constant. For some fixed $K>0$ we may then arrange
\[
 R^2=\mathrm{e}^{-K}|z_1|^2,\qquad \tau=L+K
\]
near $D$, and $\tau\geq1$ on $M$. For both profiles, $E_\tau$ is exactly the tangent space of the remaining torus factors near $D$, so $\det_\theta B_\tau=1$.

\Needspace{8\baselineskip}
\begin{example}\label{ex:critical-complete-divisor}\index{Metric completeness!critical logarithmic exponent}
Fix $\lambda>0$ and $0<A<1/\lambda$, and take $u_{\mathrm c}=\pi_1^*v_{\mathrm c}$ with the background of \cref{lem:torus-logarithmic-full-mass}. This gives a complete metric on $M$ solving
\[
 \mathrm{MA}_\theta(u_{\mathrm c})
 =\mathrm{e}^{\lambda u_{\mathrm c}}f_{\mathrm c}\theta^n,
 \qquad f_{\mathrm c}\in C^\infty(M)\cap L^1(X,\theta^n),\quad f_{\mathrm c}>0,
\]
where the right side is extended by zero on $D$. The three exponents in \eqref{eq:divisorial-three-bounds} can be taken to be
\[
 q=2-\lambda A>1,\qquad a=\lambda A,\qquad b=0,
\]
so their sum equals the critical exponent $2$.

Indeed, define $f_{\mathrm c}=\mathrm{e}^{-\lambda u_{\mathrm c}}\det_\theta(\theta+\ddc u_{\mathrm c})$ on $M$. On the inner disk the product structure gives
\[
 f_{\mathrm c}=L^{\lambda A}
       +\frac{A}{C|z_1|^2L^{2-\lambda A}}.
\]
The substitution $L=-\log r^2$ in the radial integral shows that the second term is integrable precisely when $\lambda A<1$. The first term is integrable as well, and the density is smooth elsewhere. The identity of measures across $D$ follows from the preceding lemma. Since $\tau=L+K$, all three bounds hold with the stated exponents, and completeness follows from \cref{cor:divisorial-exponential-completeness}. Along a radial normal segment the divergence of the length at the critical exponent is also visible directly:
\[
 \int_0^{r_1}\frac{\sqrt A}{r(-\log r^2)}\,\mathrm{d}r=\infty.
\]
For arbitrary paths, however, completeness requires the corollary, not just this family of radial paths.
\end{example}

\Needspace{9\baselineskip}
\begin{example}\label{ex:supercritical-incomplete-divisor}
Fix $\lambda,\eta,\beta>0$ and take $u_\beta=\pi_1^*v_\beta$ with the background of \cref{lem:torus-logarithmic-full-mass}. This is a continuous bounded solution of full mass with minimal ordinary singularities, and its metric on $M$ is incomplete, although
\[
 \mathrm{MA}_\theta(u_\beta)=\mathrm{e}^{\lambda u_\beta}f_\beta\theta^n,
 \qquad f_\beta\in C^\infty(M)\cap L^1(X,\theta^n),\quad f_\beta>0.
\]
To see this, put $K_\beta=\eta\beta(\beta+1)$ and define the density by the equation on $M$. Near $D$ it is
\[
 f_\beta=\mathrm{e}^{-\lambda\eta L^{-\beta}}
       \left(1+\frac{K_\beta}{C|z_1|^2L^{\beta+2}}\right).
\]
It is integrable, and the global identity follows as before. Fix the other coordinates and let $z_1=r\downarrow0$ along a radial segment starting at $r_0$, with $0<r_0<r_1$. The length of this segment is at most
\[
 \sqrt C\,r_0+\sqrt{K_\beta}
   \int_0^{r_0}\frac{\mathrm{d}r}{r(-\log r^2)^{1+\beta/2}}
 =\sqrt C\,r_0+\frac{\sqrt{K_\beta}}{\beta}
       (-\log r_0^2)^{-\beta/2}<\infty.
\]
Points along this segment approaching $D$ form a Cauchy sequence with no limit in $M$.

Here the smallest admissible exponents in \eqref{eq:divisorial-three-bounds} are $q=\beta+2$, $a=0$ and $b=0$: the density is asymptotic to a positive constant times $R^{-2}\tau^{-(\beta+2)}$, the potential is bounded, and the tangential determinant is $1$. Their sum is $2+\beta$. Since $\beta>0$ is arbitrary, the threshold $2$ in the criterion cannot be replaced by any larger universal number, even in a K\"ahler class with a bounded full-mass solution.
\end{example}

In these examples one may write $f=F/R^2$ with $F=R^2f$ smooth and positive on $M$, but $F$ does not extend smoothly and positively across $D$: the logarithmic factors make it tend to zero there. Thus the examples concern integrable divisorial poles weakened by logarithmic factors. They do not decide completeness for a nonintegrable pole with smooth positive coefficient across the divisor, nor the question of compatibility with the flow discussed in \cref{sec:history-completeness}.

